\color{unimain}{\section[(обязательное) Уставки функций РЗА]{(обязательное)\\Уставки функций РЗА}\label{app:set}}
\color{black}

%%%%%%%%%%%%%%%%%%%%%%%%%%%%%%%%%%%%%%%%%%%%%%%%%%%%%%%%%% ДЗ %%%%%%%%%%%%%%%%%%%%%%%%%%%%%%%%%%%%%%%%%%%%%%%%%%%%%%%%%%%%%%%%%%
\par
Параметры для настройки ДЗ приведены в таблице \ref{dz:tbl0}.

\footnotesize
\begin{longtable}{|>{\centering\arraybackslash}m{5.3cm}|>{\centering\arraybackslash}m{3.3cm}|>{\centering\arraybackslash}m{4.2cm}|>{\centering\arraybackslash}m{1.8cm}|>{\centering\arraybackslash}m{1cm}|}
\caption{Общие параметры для настройки ДЗ\hfill\vspace{-0.5\baselineskip}}\label{dz:tbl0}\\ 
\hline
\rowcolor{unimain!40}
Параметр (Параметр на ИЧМ) & Условное обозначение на схеме & Значение/ Диапазон & Единица измерения & Шаг \\ 
\hline
\endfirsthead
\caption*{\hspace{3pt}\emph{Продолжение таблицы \ref{dz:tbl0}\hfill\vspace{-0.5\baselineskip}}} \\ % сделано по ГОСТ 2.105 п.6.8.7
\hline
\rowcolor{unimain!40}
Параметр (Параметр на ИЧМ) & Условное обозначение на схеме & Значение/ Диапазон & Единица измерения & Шаг \\ 
%\hline
\endhead
%\hline
\endfoot
%\hline
\endlastfoot
%==+t1*LLN0|SV_LVDIS
\multicolumn{5}{|c|}{ ДЗ-фф 1 ст } \\ \hline 
\centering Ввод функции в работу (Ввод\_функции)  & \centering SGF1 & \centering Не предусмотрено \\ Предусмотрено & \centering -- & \centering \arraybackslash -- \\
\hline
\centering Угол наклона верхней стороны характеристики срабатывания для компенсации нагрузки (Фкн)  & \centering -- & \centering -45,0...0,0 & \centering град. & \centering \arraybackslash 0,1 \\
\hline
\centering Активное сопротивление срабатывания (Rср)  & \centering -- & \centering 0,01...500,00 & \centering Ом & \centering \arraybackslash 0,01 \\
\hline
\centering Реактивное сопротивление срабатывания (Xср)  & \centering -- & \centering 0,01...500,00 & \centering Ом & \centering \arraybackslash 0,01 \\
\hline
\centering Угол наклона характеристики срабатывания (Фхс)  & \centering -- & \centering 30,0...89,9 & \centering град. & \centering \arraybackslash 0,1 \\
\hline
\centering Режим направленности (Направленность)  & \centering SGF3 & \centering Ненаправленная \\ Прямонаправленная \\ Обратнонаправленная & \centering -- & \centering \arraybackslash -- \\
\hline
\centering Подхват от ненаправленного пуска или отдельной ступени (Подхват)  & \centering SGF9 & \centering Не предусмотрено \\ Предусмотрено & \centering -- & \centering \arraybackslash -- \\
\hline
\centering Режим контроля от БК (Реж\_БК)  & \centering SGF7 & \centering С контролем от БК-б \\ С контролем от БК-м & \centering -- & \centering \arraybackslash -- \\
\hline
\centering Режим контроля от БНН (Реж\_БНН)  & \centering SGF8 & \centering Без контроля \\ С контролем & \centering -- & \centering \arraybackslash -- \\
\hline
\centering Режим контроля от ОПО (Реж\_ОПО)  & \centering SGF6 & \centering По току \\ По току и напряжению \\ По току или по току и напряжению & \centering -- & \centering \arraybackslash -- \\
\hline
\centering Тип ПО (Контр\_ПО)  & \centering SGF5 & \centering Без контроля \\ С контролем от БК \\ С контролем от ОПО & \centering -- & \centering \arraybackslash -- \\
\hline
\centering Вывод направленности при АУ (Ненапр\_АУ)  & \centering SGF4 & \centering Не предусмотрено \\ Предусмотрено & \centering -- & \centering \arraybackslash -- \\
\hline
\centering Выдержка времени срабатывания (Tср)  & \centering T1 & \centering 0...10000 & \centering мс & \centering \arraybackslash 10 \\
\hline
\multicolumn{5}{|c|}{ ДЗ-фф 2 ст } \\ \hline 
\centering Ввод функции в работу (Ввод\_функции)  & \centering SGF1 & \centering Не предусмотрено \\ Предусмотрено & \centering -- & \centering \arraybackslash -- \\
\hline
\centering Угол наклона верхней стороны характеристики срабатывания для компенсации нагрузки (Фкн)  & \centering -- & \centering -45,0...0,0 & \centering град. & \centering \arraybackslash 0,1 \\
\hline
\centering Активное сопротивление срабатывания (Rср)  & \centering -- & \centering 0,01...500,00 & \centering Ом & \centering \arraybackslash 0,01 \\
\hline
\centering Реактивное сопротивление срабатывания (Xср)  & \centering -- & \centering 0,01...500,00 & \centering Ом & \centering \arraybackslash 0,01 \\
\hline
\centering Угол наклона характеристики срабатывания (Фхс)  & \centering -- & \centering 30,0...89,9 & \centering град. & \centering \arraybackslash 0,1 \\
\hline
\centering Режим направленности (Направленность)  & \centering SGF3 & \centering Ненаправленная \\ Прямонаправленная \\ Обратнонаправленная & \centering -- & \centering \arraybackslash -- \\
\hline
\centering Режим контроля от БК (Реж\_БК)  & \centering SGF7 & \centering С контролем от БК-б \\ С контролем от БК-м & \centering -- & \centering \arraybackslash -- \\
\hline
\centering Режим контроля от БНН (Реж\_БНН)  & \centering SGF8 & \centering Без контроля \\ С контролем & \centering -- & \centering \arraybackslash -- \\
\hline
\centering Режим контроля от ОПО (Реж\_ОПО)  & \centering SGF6 & \centering По току \\ По току и напряжению \\ По току или по току и напряжению & \centering -- & \centering \arraybackslash -- \\
\hline
\centering Тип ПО (Контр\_ПО)  & \centering SGF5 & \centering Без контроля \\ С контролем от БК \\ С контролем от ОПО & \centering -- & \centering \arraybackslash -- \\
\hline
\centering Вывод направленности при АУ (Ненапр\_АУ)  & \centering SGF4 & \centering Не предусмотрено \\ Предусмотрено & \centering -- & \centering \arraybackslash -- \\
\hline
\centering Выдержка времени срабатывания (Tср)  & \centering T1 & \centering 0...10000 & \centering мс & \centering \arraybackslash 10 \\
\hline
\multicolumn{5}{|c|}{ ДЗ-фф 3 ст } \\ \hline 
\centering Ввод функции в работу (Ввод\_функции)  & \centering SGF1 & \centering Не предусмотрено \\ Предусмотрено & \centering -- & \centering \arraybackslash -- \\
\hline
\centering Угол наклона верхней стороны характеристики срабатывания для компенсации нагрузки (Фкн)  & \centering -- & \centering -45,0...0,0 & \centering град. & \centering \arraybackslash 0,1 \\
\hline
\centering Активное сопротивление срабатывания (Rср)  & \centering -- & \centering 0,01...500,00 & \centering Ом & \centering \arraybackslash 0,01 \\
\hline
\centering Реактивное сопротивление срабатывания (Xср)  & \centering -- & \centering 0,01...500,00 & \centering Ом & \centering \arraybackslash 0,01 \\
\hline
\centering Угол наклона характеристики срабатывания (Фхс)  & \centering -- & \centering 30,0...89,9 & \centering град. & \centering \arraybackslash 0,1 \\
\hline
\centering Режим направленности (Направленность)  & \centering SGF3 & \centering Ненаправленная \\ Прямонаправленная \\ Обратнонаправленная & \centering -- & \centering \arraybackslash -- \\
\hline
\centering Подхват от ненаправленного пуска или отдельной ступени (Подхват)  & \centering SGF9 & \centering Не предусмотрено \\ Предусмотрено & \centering -- & \centering \arraybackslash -- \\
\hline
\centering Режим контроля от БК (Реж\_БК)  & \centering SGF7 & \centering С контролем от БК-б \\ С контролем от БК-м & \centering -- & \centering \arraybackslash -- \\
\hline
\centering Режим контроля от БНН (Реж\_БНН)  & \centering SGF8 & \centering Без контроля \\ С контролем & \centering -- & \centering \arraybackslash -- \\
\hline
\centering Режим контроля от ОПО (Реж\_ОПО)  & \centering SGF6 & \centering По току \\ По току и напряжению \\ По току или по току и напряжению & \centering -- & \centering \arraybackslash -- \\
\hline
\centering Тип ПО (Контр\_ПО)  & \centering SGF5 & \centering Без контроля \\ С контролем от БК \\ С контролем от ОПО & \centering -- & \centering \arraybackslash -- \\
\hline
\centering Вывод направленности при АУ (Ненапр\_АУ)  & \centering SGF4 & \centering Не предусмотрено \\ Предусмотрено & \centering -- & \centering \arraybackslash -- \\
\hline
\centering Выдержка времени срабатывания (Tср)  & \centering T1 & \centering 0...10000 & \centering мс & \centering \arraybackslash 10 \\
\hline
\multicolumn{5}{|c|}{ ДЗ-фф 4 ст } \\ \hline 
\centering Ввод функции в работу (Ввод\_функции)  & \centering SGF1 & \centering Не предусмотрено \\ Предусмотрено & \centering -- & \centering \arraybackslash -- \\
\hline
\centering Угол наклона верхней стороны характеристики срабатывания для компенсации нагрузки (Фкн)  & \centering -- & \centering -45,0...0,0 & \centering град. & \centering \arraybackslash 0,1 \\
\hline
\centering Активное сопротивление срабатывания (Rср)  & \centering -- & \centering 0,01...500,00 & \centering Ом & \centering \arraybackslash 0,01 \\
\hline
\centering Реактивное сопротивление срабатывания (Xср)  & \centering -- & \centering 0,01...500,00 & \centering Ом & \centering \arraybackslash 0,01 \\
\hline
\centering Угол наклона характеристики срабатывания (Фхс)  & \centering -- & \centering 30,0...89,9 & \centering град. & \centering \arraybackslash 0,1 \\
\hline
\centering Режим направленности (Направленность)  & \centering SGF3 & \centering Ненаправленная \\ Прямонаправленная \\ Обратнонаправленная & \centering -- & \centering \arraybackslash -- \\
\hline
\centering Режим контроля от БК (Реж\_БК)  & \centering SGF7 & \centering С контролем от БК-б \\ С контролем от БК-м & \centering -- & \centering \arraybackslash -- \\
\hline
\centering Режим контроля от БНН (Реж\_БНН)  & \centering SGF8 & \centering Без контроля \\ С контролем & \centering -- & \centering \arraybackslash -- \\
\hline
\centering Режим контроля от ОПО (Реж\_ОПО)  & \centering SGF6 & \centering По току \\ По току и напряжению \\ По току или по току и напряжению & \centering -- & \centering \arraybackslash -- \\
\hline
\centering Тип ПО (Контр\_ПО)  & \centering SGF5 & \centering Без контроля \\ С контролем от БК \\ С контролем от ОПО & \centering -- & \centering \arraybackslash -- \\
\hline
\centering Вывод направленности при АУ (Ненапр\_АУ)  & \centering SGF4 & \centering Не предусмотрено \\ Предусмотрено & \centering -- & \centering \arraybackslash -- \\
\hline
\centering Выдержка времени срабатывания (Tср)  & \centering T1 & \centering 0...10000 & \centering мс & \centering \arraybackslash 10 \\
\hline
\multicolumn{5}{|c|}{ АУ ДЗ } \\ \hline 
\centering Выбор ускоряемой ступени (Выбор\_ст)  & \centering SGF1 & \centering 1 ступень \\ 2 ступень \\ 3 ступень \\ 4 ступень \\ 1 или 3 ступень \\ 2 или 4 ступень & \centering -- & \centering \arraybackslash -- \\
\hline
\multicolumn{5}{|c|}{ ОПО } \\ \hline 
\centering Ток срабатывания ПО ДЗ по току (Iср)  & \centering I\verb|>>| & \centering 0,10...5,00 & \centering о.е. & \centering \arraybackslash 0,01 \\
\hline
\centering Ток срабатывания ПО ДЗ по току и напряжению (IсрU)  & \centering I> & \centering 0,10...5,00 & \centering о.е. & \centering \arraybackslash 0,01 \\
\hline
\centering Напряжение срабатывания ПО по току и напряжению (Uср)  & \centering U< & \centering 2,0...100,0 & \centering В & \centering \arraybackslash 0,1 \\
\hline
\centering Режим контроля от БНН (Реж\_БНН)  & \centering SGF1 & \centering Не предусмотрено \\ Предусмотрено & \centering -- & \centering \arraybackslash -- \\
\hline
\multicolumn{5}{|c|}{ ОН } \\ \hline 
\centering Угол направленности во II квадранте (Фнапр\_II)  & \centering φ2 & \centering 90,0...130,0 & \centering град. & \centering \arraybackslash 0,1 \\
\hline
\centering Угол направленности в IV квадранте (Фнапр\_IV)  & \centering φ1 & \centering -40,0...0,0 & \centering град. & \centering \arraybackslash 0,1 \\
\hline
\multicolumn{5}{|c|}{ ОУ ДЗ } \\ \hline 
\centering Выбор оперативно ускоряемой ступени (Выбор\_ст\_ОУ)  & \centering SGF1 & \centering 1 ступень \\ 2 ступень \\ 3 ступень \\ 4 ступень \\ 1 или 3 ступень \\ 2 или 4 ступень & \centering -- & \centering \arraybackslash -- \\
\hline
\centering Выдержка времени срабатывания оперативного ускорения (T\_ОУ)  & \centering T1 & \centering 0...3000 & \centering мс & \centering \arraybackslash 10 \\
\hline
\multicolumn{5}{|c|}{ ДЗ } \\ \hline 
\centering Учет нагрузочного режима в характеристике срабатывания (Вырез\_нагрузки)  & \centering SGF2 & \centering Не предусмотрено \\ Предусмотрено & \centering -- & \centering \arraybackslash -- \\
\hline
\centering Активное сопротивление для отстройки от режима максимальной нагрузки (Rнг)  & \centering -- & \centering 0,01...500,00 & \centering Ом & \centering \arraybackslash 0,01 \\
\hline
\centering Угол наклона для отстройки от режима максимальной нагрузки (Фнг)  & \centering -- & \centering 5,0...60,0 & \centering град. & \centering \arraybackslash 0,1 \\
\hline
%===t1
\end{longtable}
\normalsize



%%%%%%%%%%%%%%%%%%%%%%%%%%%%%%%%%%%%%%%%%%%%%%%%%%%%%%%%%% БК  %%%%%%%%%%%%%%%%%%%%%%%%%%%%%%%%%%%%%%%%%%%%%%%%%%%%%%%%%%%%%%%%%%
\par
В таблице \ref{bk:tbl1} приведены параметры, необходимые для настройки БК.

\footnotesize
\begin{longtable}{|>{\centering\arraybackslash}m{5.3cm}|>{\centering\arraybackslash}m{3.3cm}|>{\centering\arraybackslash}m{4.2cm}|>{\centering\arraybackslash}m{1.8cm}|>{\centering\arraybackslash}m{1cm}|}
\caption{Параметры для настройки БК\hfill\vspace{-0.5\baselineskip}}\label{bk:tbl1}\\ 
\hline
\rowcolor{gray!20}
Параметр (Параметр на ИЧМ) & Условное обозначение на схеме & Значение/ Диапазон & Единица измерения & Шаг \\ 
\hline
\endfirsthead
\caption*{\hspace{3pt}\emph{Продолжение таблицы \ref{bk:tbl1}\hfill\vspace{-0.5\baselineskip}}} \\ % сделано по ГОСТ 2.105 п.6.8.7
\hline
\rowcolor{gray!20}
Параметр (Параметр на ИЧМ) & Условное обозначение на схеме & Значение/ Диапазон & Единица измерения & Шаг \\ 
%\hline
\endhead
%\hline
\endfoot
%\hline
\endlastfoot
%==+t1*RDIS1|LVPSB
\centering Ввод функции в работу (Ввод\_функции)  & \centering SGF1 & \centering Не предусмотрено \\ Предусмотрено & \centering -- & \centering \arraybackslash -- \\
\hline
\centering Приращение тока ПП чувствительного органа (dI1\_чув)  & \centering dI1/dt> & \centering 0,08...3,00 & \centering о.е. & \centering \arraybackslash 0,01 \\
\hline
\centering Приращение тока ПП грубого органа (dI1\_груб)  & \centering dI1/dt\verb|>>| & \centering 0,12...5,00 & \centering о.е. & \centering \arraybackslash 0,01 \\
\hline
\centering Приращение тока ОП чувствительного органа (dI2\_чув)  & \centering dI2/dt> & \centering 0,04...1,50 & \centering о.е. & \centering \arraybackslash 0,01 \\
\hline
\centering Приращение тока ОП грубого органа (dI2\_груб)  & \centering dI2/dt\verb|>>| & \centering 0,06...2,50 & \centering о.е. & \centering \arraybackslash 0,01 \\
\hline
\centering Время ввода быстродействующих ступеней ДЗ от чувствительных органов (Тб\_чув)  & \centering T1 & \centering 200...1000 & \centering мс & \centering \arraybackslash 10 \\
\hline
\centering Время ввода быстродействующих ступеней ДЗ от грубых органов (Тб\_груб)  & \centering T2 & \centering 200...1000 & \centering мс & \centering \arraybackslash 10 \\
\hline
\centering Время ввода медленнодействующих ступеней ДЗ (Тм)  & \centering T3 & \centering 1000...15000 & \centering мс & \centering \arraybackslash 10 \\
\hline
\centering Ускоренный возврат логики БК при отключении выключателя (Ускор\_возврат)  & \centering SGF2 & \centering Не предусмотрено \\ Предусмотрено & \centering -- & \centering \arraybackslash -- \\
\hline
%===t1
\end{longtable}
\normalsize



%%%%%%%%%%%%%%%%%%%%%%%%%%%%%%%%%%%%%%%%%%%%%%%%%%%%%%%%%% АУ %%%%%%%%%%%%%%%%%%%%%%%%%%%%%%%%%%%%%%%%%%%%%%%%%%%%%%%%%%%%%%%%%%
\par
В таблице \ref{au:tbl1} приведены параметры, необходимые для настройки АУ.

\footnotesize
\begin{longtable}{|>{\centering\arraybackslash}m{5.3cm}|>{\centering\arraybackslash}m{3.3cm}|>{\centering\arraybackslash}m{4.2cm}|>{\centering\arraybackslash}m{1.8cm}|>{\centering\arraybackslash}m{1cm}|}
\caption{Параметры для настройки АУ\hfill\vspace{-0.5\baselineskip}}\label{au:tbl1}\\ 
\hline
\rowcolor{gray!20}
Параметр (Параметр на ИЧМ) & Условное обозначение на схеме & Значение/ Диапазон & Единица измерения & Шаг \\ 
\hline
\endfirsthead
\caption*{\hspace{3pt}\emph{Продолжение таблицы \ref{au:tbl1}\hfill\vspace{-0.5\baselineskip}}} \\ % сделано по ГОСТ 2.105 п.6.8.7
\hline
\rowcolor{gray!20}
Параметр (Параметр на ИЧМ) & Условное обозначение на схеме & Значение/ Диапазон & Единица измерения & Шаг \\ 
%\hline
\endhead
%\hline
\endfoot
%\hline
\endlastfoot
%==+t1*PSOF1|LVLINAUA
\centering Ввод функции в работу (Ввод\_функции)  & \centering SGF1 & \centering Не предусмотрено \\ Предусмотрено & \centering -- & \centering \arraybackslash -- \\
\hline
\centering Автоматическое ускорение МТЗ (АУ\_МТЗ)  & \centering SGF2 & \centering Не предусмотрено \\ Предусмотрено & \centering -- & \centering \arraybackslash -- \\
\hline
\centering Автоматическое ускорение ДЗ (АУ\_ДЗ)  & \centering SGF3 & \centering Не предусмотрено \\ Предусмотрено & \centering -- & \centering \arraybackslash -- \\
\hline
\centering Автоматическое ускорение ТЗНП (АУ\_ТЗНП)  & \centering SGF4 & \centering Не предусмотрено \\ Предусмотрено & \centering -- & \centering \arraybackslash -- \\
\hline
\centering Выдержка времени срабатывания ускоряемой ступени МТЗ (Тср\_МТЗ)  & \centering T1 & \centering 0...3000 & \centering мс & \centering \arraybackslash 10 \\
\hline
\centering Выдержка времени срабатывания ускоряемой ступени ДЗ (Тср\_ДЗ)  & \centering T2 & \centering 0...3000 & \centering мс & \centering \arraybackslash 10 \\
\hline
\centering Выдержка времени срабатывания ускоряемой ступени ТЗНП (Тср\_ТЗНП)  & \centering T3 & \centering 0...3000 & \centering мс & \centering \arraybackslash 10 \\
\hline
\centering Время ввода ускорения при включении выключателя (Твв)  & \centering T4 & \centering 0...5000 & \centering мс & \centering \arraybackslash 10 \\
\hline
%===t1
\end{longtable}
\normalsize


%%%%%%%%%%%%%%%%%%%%%%%%%%%%%%%%%%%%%%%%%%%%%%%%%%%%%%%%%% МТЗ %%%%%%%%%%%%%%%%%%%%%%%%%%%%%%%%%%%%%%%%%%%%%%%%%%%%%%%%%%%%%%%%%% \footnotesize
\par
В таблице \ref{mtz:tbl1} приведены параметры, необходимые для настройки МТЗ.

\footnotesize
\begin{longtable}{|>{\centering\arraybackslash}m{5.3cm}|>{\centering\arraybackslash}m{3.3cm}|>{\centering\arraybackslash}m{4.2cm}|>{\centering\arraybackslash}m{1.8cm}|>{\centering\arraybackslash}m{1cm}|}
\caption{Параметры для настройки ступени МТЗ\hfill\vspace{-0.5\baselineskip}}\label{mtz:tbl1}\\ 
\hline
\rowcolor{unimain!40}
Параметр (Параметр на ИЧМ) & Условное обозначение на схеме & Значение/ Диапазон & Единица измерения & Шаг \\ 
\hline
\endfirsthead
\caption*{\hspace{3pt}\emph{Продолжение таблицы \ref{mtz:tbl1}\hfill\vspace{-0.5\baselineskip}}} \\ % сделано по ГОСТ 2.105 п.6.8.7
\hline
\rowcolor{unimain!40}
Параметр (Параметр на ИЧМ) & Условное обозначение на схеме & Значение/ Диапазон & Единица измерения & Шаг \\ 
%\hline
\endhead
%\hline
\endfoot
%\hline
\endlastfoot
%==+t1*PTOC1|SV_LVTUVTOC
\multicolumn{5}{|c|}{ МТЗ 1 ст } \\ \hline 
\centering Ввод функции в работу (Ввод\_функции)  & \centering SGF1 & \centering Не предусмотрено \\ Предусмотрено & \centering -- & \centering \arraybackslash -- \\
\hline
\centering Ток срабатывания (Iср)  & \centering I> & \centering 0,10...30,00 & \centering о.е. & \centering \arraybackslash 0,01 \\
\hline
\centering Ток срабатывания грубого органа в режиме управляющего напряжения (Iср\_груб.)  & \centering -- & \centering 0,10...30,00 & \centering о.е. & \centering \arraybackslash 0,01 \\
\hline
\centering Ток срабатывания функции включения под нагрузку (Iфвн)  & \centering -- & \centering 0,10...30,00 & \centering о.е. & \centering \arraybackslash 0,01 \\
\hline
\centering Ввод функции включения на "холодную" нагрузку (Ввод\_ФВН)  & \centering SGF2 & \centering Не предусмотрено \\ Предусмотрено & \centering -- & \centering \arraybackslash -- \\
\hline
\centering Время отключенного состояния выключателя (Тзагруб.)  & \centering T2 & \centering 0...20000 & \centering мс & \centering \arraybackslash 10 \\
\hline
\centering Время ввода загрубления (Тсброс.загруб.)  & \centering T3 & \centering 0...20000 & \centering мс & \centering \arraybackslash 10 \\
\hline
\centering Тип пуска по напряжению (Тип\_КПОН)  & \centering SGF6 & \centering Управляющее напряжение \\ Вольтметровая блокировка & \centering -- & \centering \arraybackslash -- \\
\hline
\centering Режим контроля от БНТ (Реж\_БНТ)  & \centering SGF7 & \centering Не предусмотрено \\ Предусмотрено & \centering -- & \centering \arraybackslash -- \\
\hline
\centering Режим направленности (Направленность)  & \centering SGF3 & \centering Ненаправленная \\ Прямонаправленная \\ Обратнонаправленная & \centering -- & \centering \arraybackslash -- \\
\hline
\centering Режим ОНМ при неисправности ЦН (Реж\_БНН\_ОНМ)  & \centering SGF4 & \centering Блокировка \\ Деблокировка & \centering -- & \centering \arraybackslash -- \\
\hline
\centering Вывод направленности при АУ (Ненапр\_АУ)  & \centering SGF5 & \centering Не предусмотрено \\ Предусмотрено & \centering -- & \centering \arraybackslash -- \\
\hline
\centering Режим КПОН при неисправности ЦН (Реж\_БНН\_КПОН)  & \centering SGF8 & \centering Деблокировка \\ Блокировка & \centering -- & \centering \arraybackslash -- \\
\hline
\centering Режим контроля от КПОН (Реж\_КПОН)  & \centering SGF9 & \centering Не предусмотрено \\ Предусмотрено & \centering -- & \centering \arraybackslash -- \\
\hline
\centering Характеристика срабатывания (Характеристика)  & \centering -- & \centering Кривая F \\ Кривая E \\ Кривая D \\ Кривая B \\ Кривая A \\ Кривая C \\ Независимая \\ Кривая RXIDG & \centering -- & \centering \arraybackslash -- \\
\hline
\centering Коэффициент регулировки времени срабатывания зависимой характеристики (К)  & \centering -- & \centering 0,05...15,00 & \centering -- & \centering \arraybackslash 0,01 \\
\hline
\centering Независимая выдержка времени срабатывания (Tср)  & \centering T1 & \centering 0...100000 & \centering мс & \centering \arraybackslash 10 \\
\hline
\multicolumn{5}{|c|}{ МТЗ 2 ст } \\ \hline 
\centering Ввод функции в работу (Ввод\_функции)  & \centering SGF1 & \centering Не предусмотрено \\ Предусмотрено & \centering -- & \centering \arraybackslash -- \\
\hline
\centering Ток срабатывания (Iср)  & \centering I> & \centering 0,10...30,00 & \centering о.е. & \centering \arraybackslash 0,01 \\
\hline
\centering Ток срабатывания грубого органа в режиме управляющего напряжения (Iср\_груб.)  & \centering -- & \centering 0,10...30,00 & \centering о.е. & \centering \arraybackslash 0,01 \\
\hline
\centering Ток срабатывания функции включения под нагрузку (Iфвн)  & \centering -- & \centering 0,10...30,00 & \centering о.е. & \centering \arraybackslash 0,01 \\
\hline
\centering Ввод функции включения на "холодную" нагрузку (Ввод\_ФВН)  & \centering SGF2 & \centering Не предусмотрено \\ Предусмотрено & \centering -- & \centering \arraybackslash -- \\
\hline
\centering Время отключенного состояния выключателя (Тзагруб.)  & \centering T2 & \centering 0...20000 & \centering мс & \centering \arraybackslash 10 \\
\hline
\centering Время ввода загрубления (Тсброс.загруб.)  & \centering T3 & \centering 0...20000 & \centering мс & \centering \arraybackslash 10 \\
\hline
\centering Тип пуска по напряжению (Тип\_КПОН)  & \centering SGF6 & \centering Управляющее напряжение \\ Вольтметровая блокировка & \centering -- & \centering \arraybackslash -- \\
\hline
\centering Режим контроля от БНТ (Реж\_БНТ)  & \centering SGF7 & \centering Не предусмотрено \\ Предусмотрено & \centering -- & \centering \arraybackslash -- \\
\hline
\centering Режим направленности (Направленность)  & \centering SGF3 & \centering Ненаправленная \\ Прямонаправленная \\ Обратнонаправленная & \centering -- & \centering \arraybackslash -- \\
\hline
\centering Режим ОНМ при неисправности ЦН (Реж\_БНН\_ОНМ)  & \centering SGF4 & \centering Блокировка \\ Деблокировка & \centering -- & \centering \arraybackslash -- \\
\hline
\centering Вывод направленности при АУ (Ненапр\_АУ)  & \centering SGF5 & \centering Не предусмотрено \\ Предусмотрено & \centering -- & \centering \arraybackslash -- \\
\hline
\centering Режим КПОН при неисправности ЦН (Реж\_БНН\_КПОН)  & \centering SGF8 & \centering Деблокировка \\ Блокировка & \centering -- & \centering \arraybackslash -- \\
\hline
\centering Режим контроля от КПОН (Реж\_КПОН)  & \centering SGF9 & \centering Не предусмотрено \\ Предусмотрено & \centering -- & \centering \arraybackslash -- \\
\hline
\centering Характеристика срабатывания (Характеристика)  & \centering -- & \centering Кривая F \\ Кривая E \\ Кривая D \\ Кривая B \\ Кривая A \\ Кривая C \\ Независимая \\ Кривая RXIDG & \centering -- & \centering \arraybackslash -- \\
\hline
\centering Коэффициент регулировки времени срабатывания зависимой характеристики (К)  & \centering -- & \centering 0,05...15,00 & \centering -- & \centering \arraybackslash 0,01 \\
\hline
\centering Независимая выдержка времени срабатывания (Tср)  & \centering T1 & \centering 0...100000 & \centering мс & \centering \arraybackslash 10 \\
\hline
\multicolumn{5}{|c|}{ МТЗ 3 ст } \\ \hline 
\centering Ввод функции в работу (Ввод\_функции)  & \centering SGF1 & \centering Не предусмотрено \\ Предусмотрено & \centering -- & \centering \arraybackslash -- \\
\hline
\centering Ток срабатывания (Iср)  & \centering I> & \centering 0,10...30,00 & \centering о.е. & \centering \arraybackslash 0,01 \\
\hline
\centering Ток срабатывания грубого органа в режиме управляющего напряжения (Iср\_груб.)  & \centering -- & \centering 0,10...30,00 & \centering о.е. & \centering \arraybackslash 0,01 \\
\hline
\centering Ток срабатывания функции включения под нагрузку (Iфвн)  & \centering -- & \centering 0,10...30,00 & \centering о.е. & \centering \arraybackslash 0,01 \\
\hline
\centering Ввод функции включения на "холодную" нагрузку (Ввод\_ФВН)  & \centering SGF2 & \centering Не предусмотрено \\ Предусмотрено & \centering -- & \centering \arraybackslash -- \\
\hline
\centering Время отключенного состояния выключателя (Тзагруб.)  & \centering T2 & \centering 0...20000 & \centering мс & \centering \arraybackslash 10 \\
\hline
\centering Время ввода загрубления (Тсброс.загруб.)  & \centering T3 & \centering 0...20000 & \centering мс & \centering \arraybackslash 10 \\
\hline
\centering Тип пуска по напряжению (Тип\_КПОН)  & \centering SGF6 & \centering Управляющее напряжение \\ Вольтметровая блокировка & \centering -- & \centering \arraybackslash -- \\
\hline
\centering Режим контроля от БНТ (Реж\_БНТ)  & \centering SGF7 & \centering Не предусмотрено \\ Предусмотрено & \centering -- & \centering \arraybackslash -- \\
\hline
\centering Режим направленности (Направленность)  & \centering SGF3 & \centering Ненаправленная \\ Прямонаправленная \\ Обратнонаправленная & \centering -- & \centering \arraybackslash -- \\
\hline
\centering Режим ОНМ при неисправности ЦН (Реж\_БНН\_ОНМ)  & \centering SGF4 & \centering Блокировка \\ Деблокировка & \centering -- & \centering \arraybackslash -- \\
\hline
\centering Вывод направленности при АУ (Ненапр\_АУ)  & \centering SGF5 & \centering Не предусмотрено \\ Предусмотрено & \centering -- & \centering \arraybackslash -- \\
\hline
\centering Режим КПОН при неисправности ЦН (Реж\_БНН\_КПОН)  & \centering SGF8 & \centering Деблокировка \\ Блокировка & \centering -- & \centering \arraybackslash -- \\
\hline
\centering Режим контроля от КПОН (Реж\_КПОН)  & \centering SGF9 & \centering Не предусмотрено \\ Предусмотрено & \centering -- & \centering \arraybackslash -- \\
\hline
\centering Характеристика срабатывания (Характеристика)  & \centering -- & \centering Кривая F \\ Кривая E \\ Кривая D \\ Кривая B \\ Кривая A \\ Кривая C \\ Независимая \\ Кривая RXIDG & \centering -- & \centering \arraybackslash -- \\
\hline
\centering Коэффициент регулировки времени срабатывания зависимой характеристики (К)  & \centering -- & \centering 0,05...15,00 & \centering -- & \centering \arraybackslash 0,01 \\
\hline
\centering Независимая выдержка времени срабатывания (Tср)  & \centering T1 & \centering 0...100000 & \centering мс & \centering \arraybackslash 10 \\
\hline
\multicolumn{5}{|c|}{ МТЗ 4 ст } \\ \hline 
\centering Ввод функции в работу (Ввод\_функции)  & \centering SGF1 & \centering Не предусмотрено \\ Предусмотрено & \centering -- & \centering \arraybackslash -- \\
\hline
\centering Ток срабатывания (Iср)  & \centering I> & \centering 0,10...30,00 & \centering о.е. & \centering \arraybackslash 0,01 \\
\hline
\centering Ток срабатывания грубого органа в режиме управляющего напряжения (Iср\_груб.)  & \centering -- & \centering 0,10...30,00 & \centering о.е. & \centering \arraybackslash 0,01 \\
\hline
\centering Ток срабатывания функции включения под нагрузку (Iфвн)  & \centering -- & \centering 0,10...30,00 & \centering о.е. & \centering \arraybackslash 0,01 \\
\hline
\centering Ввод функции включения на "холодную" нагрузку (Ввод\_ФВН)  & \centering SGF2 & \centering Не предусмотрено \\ Предусмотрено & \centering -- & \centering \arraybackslash -- \\
\hline
\centering Время отключенного состояния выключателя (Тзагруб.)  & \centering T2 & \centering 0...20000 & \centering мс & \centering \arraybackslash 10 \\
\hline
\centering Время ввода загрубления (Тсброс.загруб.)  & \centering T3 & \centering 0...20000 & \centering мс & \centering \arraybackslash 10 \\
\hline
\centering Тип пуска по напряжению (Тип\_КПОН)  & \centering SGF6 & \centering Управляющее напряжение \\ Вольтметровая блокировка & \centering -- & \centering \arraybackslash -- \\
\hline
\centering Режим контроля от БНТ (Реж\_БНТ)  & \centering SGF7 & \centering Не предусмотрено \\ Предусмотрено & \centering -- & \centering \arraybackslash -- \\
\hline
\centering Режим направленности (Направленность)  & \centering SGF3 & \centering Ненаправленная \\ Прямонаправленная \\ Обратнонаправленная & \centering -- & \centering \arraybackslash -- \\
\hline
\centering Режим ОНМ при неисправности ЦН (Реж\_БНН\_ОНМ)  & \centering SGF4 & \centering Блокировка \\ Деблокировка & \centering -- & \centering \arraybackslash -- \\
\hline
\centering Вывод направленности при АУ (Ненапр\_АУ)  & \centering SGF5 & \centering Не предусмотрено \\ Предусмотрено & \centering -- & \centering \arraybackslash -- \\
\hline
\centering Режим КПОН при неисправности ЦН (Реж\_БНН\_КПОН)  & \centering SGF8 & \centering Деблокировка \\ Блокировка & \centering -- & \centering \arraybackslash -- \\
\hline
\centering Режим контроля от КПОН (Реж\_КПОН)  & \centering SGF9 & \centering Не предусмотрено \\ Предусмотрено & \centering -- & \centering \arraybackslash -- \\
\hline
\centering Характеристика срабатывания (Характеристика)  & \centering -- & \centering Кривая F \\ Кривая E \\ Кривая D \\ Кривая B \\ Кривая A \\ Кривая C \\ Независимая \\ Кривая RXIDG & \centering -- & \centering \arraybackslash -- \\
\hline
\centering Коэффициент регулировки времени срабатывания зависимой характеристики (К)  & \centering -- & \centering 0,05...15,00 & \centering -- & \centering \arraybackslash 0,01 \\
\hline
\centering Независимая выдержка времени срабатывания (Tср)  & \centering T1 & \centering 0...100000 & \centering мс & \centering \arraybackslash 10 \\
\hline
\multicolumn{5}{|c|}{ КПОН } \\ \hline 
\centering Напряжение срабатывания (Uср)  & \centering U< & \centering 2,0...100,0 & \centering В & \centering \arraybackslash 0,1 \\
\hline
\centering Напряжение срабатывания ОП (U2ср)  & \centering U2> & \centering 2,0...50,0 & \centering В & \centering \arraybackslash 0,1 \\
\hline
\centering Режим пуска (Реж\_пуска)  & \centering SGF1 & \centering По Uмин \\ Комбинированный \\ Внешний & \centering -- & \centering \arraybackslash -- \\
\hline
\multicolumn{5}{|c|}{ ОНМ } \\ \hline 
\centering Угол максимальной чувствительности (Фмч)  & \centering φмч & \centering -179,0...180,0 & \centering град. & \centering \arraybackslash 0,1 \\
\hline
\multicolumn{5}{|c|}{ БНТ } \\ \hline 
\centering Отношение тока второй гармоники к току основной гармоники (Ih2/Ih1)  & \centering I2г/I1г> & \centering 5...100 & \centering \% & \centering \arraybackslash 1 \\
\hline
\centering Ток блокировки (Iср)  & \centering I> & \centering 0,1...15,0 & \centering о.е. & \centering \arraybackslash 0,1 \\
\hline
\centering Перекрестная блокировка (Перекрест\_блок)  & \centering SGF1 & \centering Не предусмотрено \\ Предусмотрено & \centering -- & \centering \arraybackslash -- \\
\hline
\multicolumn{5}{|c|}{ БЛЗШ } \\ \hline 
\centering Выбор ступени блокировки (БлокЛЗШ\_выбор\_ст)  & \centering SGF1 & \centering Не предусмотрено \\ 1 ступень \\ 2 ступень \\ 3 ступень \\ 4 ступень \\ 1 или 3 ступень \\ 2 или 4 ступень & \centering -- & \centering \arraybackslash -- \\
\hline
\multicolumn{5}{|c|}{ АУ МТЗ } \\ \hline 
\centering Выбор ускоряемой ступени (Выбор\_ст)  & \centering SGF1 & \centering 1 ступень \\ 2 ступень \\ 3 ступень \\ 4 ступень \\ 1 или 3 ступень \\ 2 или 4 ступень & \centering -- & \centering \arraybackslash -- \\
\hline
\multicolumn{5}{|c|}{ ОУ МТЗ } \\ \hline 
\centering Выбор оперативно ускоряемой ступени (Выбор\_ст\_ОУ)  & \centering SGF1 & \centering 1 ступень \\ 2 ступень \\ 3 ступень \\ 4 ступень \\ 1 или 3 ступень \\ 2 или 4 ступень & \centering -- & \centering \arraybackslash -- \\
\hline
\centering Выдержка времени срабатывания оперативного ускорения (T\_ОУ)  & \centering T1 & \centering 0...3000 & \centering мс & \centering \arraybackslash 10 \\
\hline
%===t1
\end{longtable}
\normalsize



%%%%%%%%%%%%%%%%%%%%%%%%%%%%%%%%%%%%%%%%%%%%%%%%%%%%%%%%%% ЛЗШ  %%%%%%%%%%%%%%%%%%%%%%%%%%%%%%%%%%%%%%%%%%%%%%%%%%%%%%%%%%%%%%%%%% \footnotesize
\par
В таблице \ref{lzsh:tbl1} приведены параметры, необходимые для настройки ЛЗШ.

\footnotesize
\begin{longtable}{|>{\centering\arraybackslash}m{5.3cm}|>{\centering\arraybackslash}m{3.3cm}|>{\centering\arraybackslash}m{4.2cm}|>{\centering\arraybackslash}m{1.8cm}|>{\centering\arraybackslash}m{1cm}|}
\caption{Параметры для настройки ЛЗШ\hfill\vspace{-0.5\baselineskip}}\label{lzsh:tbl1}\\ 
\hline
\rowcolor{gray!20}
Параметр (Параметр на ИЧМ) & Условное обозначение на схеме & Значение/ Диапазон & Единица измерения & Шаг \\ 
\hline
\endfirsthead
\caption*{\hspace{3pt}\emph{Продолжение таблицы \ref{lzsh:tbl1}\hfill\vspace{-0.5\baselineskip}}} \\ % сделано по ГОСТ 2.105 п.6.8.7
\hline
\rowcolor{gray!20}
Параметр (Параметр на ИЧМ) & Условное обозначение на схеме & Значение/ Диапазон & Единица измерения & Шаг \\ 
%\hline
\endhead
%\hline
\endfoot
%\hline
\endlastfoot
%==+t1*PTOC1|LBPTOC
\centering Ввод функции в работу (Ввод\_функции)  & \centering SGF1 & \centering Не предусмотрено \\ Предусмотрено & \centering -- & \centering \arraybackslash -- \\
\hline
\centering Ток срабатывания (Iср)  & \centering I> & \centering 0,10...30,00 & \centering о.е. & \centering \arraybackslash 0,01 \\
\hline
\centering Ток срабатывания грубого органа в режиме управляющего напряжения (Iср\_груб.)  & \centering -- & \centering 0,10...30,00 & \centering о.е. & \centering \arraybackslash 0,01 \\
\hline
\centering Тип пуска по напряжению (Тип\_КПОН)  & \centering SGF2 & \centering Управляющее напряжение \\ Вольтметровая блокировка & \centering -- & \centering \arraybackslash -- \\
\hline
\centering Режим контроля от БНТ (Реж\_БНТ)  & \centering SGF3 & \centering Не предусмотрено \\ Предусмотрено & \centering -- & \centering \arraybackslash -- \\
\hline
\centering Режим КПОН при неисправности ЦН (Реж\_БНН\_КПОН)  & \centering SGF4 & \centering Деблокировка \\ Блокировка & \centering -- & \centering \arraybackslash -- \\
\hline
\centering Режим контроля от КПОН (Реж\_КПОН)  & \centering SGF5 & \centering Не предусмотрено \\ Предусмотрено & \centering -- & \centering \arraybackslash -- \\
\hline
\centering Выдержка времени неисправности ЛЗШ (Тнеисп)  & \centering T2 & \centering 100...10000 & \centering мс & \centering \arraybackslash 10 \\
\hline
\centering Независимая выдержка времени срабатывания (Tср)  & \centering T1 & \centering 100...10000 & \centering мс & \centering \arraybackslash 10 \\
\hline
%===t1
\end{longtable}
\normalsize


%%%%%%%%%%%%%%%%%%%%%%%%%%%%%%%%%%%%%%%%%%%%%%%%%%%%%%%%%% ЗДЗ  %%%%%%%%%%%%%%%%%%%%%%%%%%%%%%%%%%%%%%%%%%%%%%%%%%%%%%%%%%%%%%%%%% \footnotesize
\par
В таблице \ref{zdz:tbl1} приведены параметры, необходимые для настройки ЗДЗ.

\footnotesize
\begin{longtable}{|>{\centering\arraybackslash}m{5.3cm}|>{\centering\arraybackslash}m{3.3cm}|>{\centering\arraybackslash}m{4.2cm}|>{\centering\arraybackslash}m{1.8cm}|>{\centering\arraybackslash}m{1cm}|}
\caption{Параметры для настройки ЗДЗ\hfill\vspace{-0.5\baselineskip}}\label{zdz:tbl1}\\ 
\hline
\rowcolor{unimain!40}
Параметр (Параметр на ИЧМ) & Условное обозначение на схеме & Значение/ Диапазон & Единица измерения & Шаг \\ 
\hline
\endfirsthead
\caption*{\hspace{3pt}\emph{Продолжение таблицы \ref{zdz:tbl1}\hfill\vspace{-0.5\baselineskip}}} \\ % сделано по ГОСТ 2.105 п.6.8.7
\hline
\rowcolor{unimain!40}
Параметр (Параметр на ИЧМ) & Условное обозначение на схеме & Значение/ Диапазон & Единица измерения & Шаг \\ 
%\hline
\endhead
%\hline
\endfoot
%\hline
\endlastfoot
%==+t1*ARCPTOC1|SV_ARCTOC
\centering Ввод функции в работу (Ввод\_функции)  & \centering SGF1 & \centering Не предусмотрено \\ Предусмотрено & \centering -- & \centering \arraybackslash -- \\
\hline
\centering Ток срабатывания (Iср)  & \centering I> & \centering 0,10...30,00 & \centering о.е. & \centering \arraybackslash 0,01 \\
\hline
\centering Напряжение срабатывания (3U0ср)  & \centering 3U0> & \centering 2,0...150,0 & \centering В & \centering \arraybackslash 0,1 \\
\hline
\centering Блокировка при неисправности ЗДЗ (Блок\_неиспр\_ЗДЗ)  & \centering SGF2 & \centering Не предусмотрено \\ Предусмотрено & \centering -- & \centering \arraybackslash -- \\
\hline
\centering Контроль тока (Контр\_тока)  & \centering SGF3 & \centering Не предусмотрено \\ Контроль I \\ Контроль I или 3U0 \\ Внешний & \centering -- & \centering \arraybackslash -- \\
\hline
\centering Режим контроля от БНН (Реж\_БНН)  & \centering SGF4 & \centering Без контроля \\ С контролем & \centering -- & \centering \arraybackslash -- \\
\hline
\centering Выдержка времени неисправности ЗДЗ (Tнеисп)  & \centering T2 & \centering 0...10000 & \centering мс & \centering \arraybackslash 10 \\
\hline
\centering Выдержка времени срабатывания (Tср)  & \centering T1 & \centering 0...10000 & \centering мс & \centering \arraybackslash 10 \\
\hline
%===t1
\end{longtable}
\normalsize


%%%%%%%%%%%%%%%%%%%%%%%%%%%%%%%%%%%%%%%%%%%%%%%%%%%%%%%%%% ТК ЗДЗ %%%%%%%%%%%%%%%%%%%%%%%%%%%%%%%%%%%%%%%%%%%%%%%%%%%%%%%%%%%%%%%%%%

В таблице \ref{tkzdz:tbl1} приведены параметры, необходимые для настройки ТК ЗДЗ.

\footnotesize
\begin{longtable}{|>{\centering\arraybackslash}m{5.3cm}|>{\centering\arraybackslash}m{3.3cm}|>{\centering\arraybackslash}m{4.2cm}|>{\centering\arraybackslash}m{1.8cm}|>{\centering\arraybackslash}m{1cm}|}
\caption{Параметры для настройки ТК ЗДЗ\hfill\vspace{-0.5\baselineskip}}\label{tkzdz:tbl1}\\ 
\hline
\rowcolor{gray!20}
Параметр (Параметр на ИЧМ) & Условное обозначение на схеме & Значение/ Диапазон & Единица измерения & Шаг \\ 
\hline
\endfirsthead
%\caption{\emph{Продолжение\hfill\vspace{-0.5\baselineskip}}} \\ % Отступ обозначения таблицы от таблицы и выравнивание надписи слева
\caption*{\hspace{3pt}\emph{Продолжение таблицы \ref{tkzdz:tbl1}\hfill\vspace{-0.5\baselineskip}}} \\ % сделано по ГОСТ 2.105 п.6.8.7
\hline
\rowcolor{gray!20}
Параметр (Параметр на ИЧМ) & Условное обозначение на схеме & Значение/ Диапазон & Единица измерения & Шаг \\ 
%\hline
\endhead
%\hline
\endfoot
%\hline
\endlastfoot
%==+t1*PTOC1|LVARCTOC
\centering Ввод функции в работу (Ввод\_функции)  & \centering SGF1 & \centering Не предусмотрено \\ Предусмотрено & \centering -- & \centering \arraybackslash -- \\
\hline
\centering Ток срабатывания (Iср)  & \centering I> & \centering 0,10...30,00 & \centering о.е. & \centering \arraybackslash 0,01 \\
\hline
%===t1
\end{longtable}
\normalsize



%%%%%%%%%%%%%%%%%%%%%%%%%%%%%%%%%%%%%%%%%%%%%%%%%%%%%%%%%% ЗОП %%%%%%%%%%%%%%%%%%%%%%%%%%%%%%%%%%%%%%%%%%%%%%%%%%%%%%%%%%%%%%%%%%
\par
В таблице \ref{zop:tbl1} приведены параметры, необходимые для настройки ЗОП.

\footnotesize
\begin{longtable}{|>{\centering\arraybackslash}m{5.3cm}|>{\centering\arraybackslash}m{3.3cm}|>{\centering\arraybackslash}m{4.2cm}|>{\centering\arraybackslash}m{1.8cm}|>{\centering\arraybackslash}m{1cm}|}
\caption{Параметры для настройки ЗОП\hfill\vspace{-0.5\baselineskip}}\label{zop:tbl1}\\ 
\hline
\rowcolor{gray!20}
Параметр (Параметр на ИЧМ) & Условное обозначение на схеме & Значение/ Диапазон & Единица измерения & Шаг \\ 
\hline
\endfirsthead
\caption*{\hspace{3pt}\emph{Продолжение таблицы \ref{zop:tbl1}\hfill\vspace{-0.5\baselineskip}}} \\ % сделано по ГОСТ 2.105 п.6.8.7
\hline
\rowcolor{gray!20}
Параметр (Параметр на ИЧМ) & Условное обозначение на схеме & Значение/ Диапазон & Единица измерения & Шаг \\ 
%\hline
\endhead
%\hline
\endfoot
%\hline
\endlastfoot
%==+t1*NSPTOC1|LVNSTOC
\centering Ввод функции в работу (Ввод\_функции)  & \centering SGF1 & \centering Не предусмотрено \\ Предусмотрено & \centering -- & \centering \arraybackslash -- \\
\hline
\centering Ток срабатывания ОП (I2ср)  & \centering I2> & \centering 0,04...4,00 & \centering о.е. & \centering \arraybackslash 0,01 \\
\hline
\centering Отношение тока ОП к току ПП (Kнесим)  & \centering I2/I1> & \centering 0,02...1,00 & \centering о.е. & \centering \arraybackslash 0,01 \\
\hline
\centering Режим работы ЗОП (Реж\_ЗОП)  & \centering SGF2 & \centering по I2 \\ по I2/I1 \\ по I2 или по I2/I1 & \centering -- & \centering \arraybackslash -- \\
\hline
\centering Выдержка времени срабатывания (Tср)  & \centering T1 & \centering 100...20000 & \centering мс & \centering \arraybackslash 10 \\
\hline
%===t1
\end{longtable}
\normalsize


%%%%%%%%%%%%%%%%%%%%%%%%%%%%%%%%%%%%%%%%%%%%%%%%%%%%%%%%%% ТЗНП %%%%%%%%%%%%%%%%%%%%%%%%%%%%%%%%%%%%%%%%%%%%%%%%%%%%%%%%%%%%%%%%%% \footnotesize
\par

В таблице \ref{tznp:tbl1} приведены параметры, необходимые для настройки ТЗНП.

\footnotesize
\begin{longtable}{|>{\centering\arraybackslash}m{5.3cm}|>{\centering\arraybackslash}m{3.3cm}|>{\centering\arraybackslash}m{4.2cm}|>{\centering\arraybackslash}m{1.8cm}|>{\centering\arraybackslash}m{1cm}|}
\caption{Параметры для настройки ступеней ТЗНП\hfill\vspace{-0.5\baselineskip}}\label{tznp:tbl1}\\ 
\hline
\rowcolor{unimain!40}
Параметр (Параметр на ИЧМ) & Условное обозначение на схеме & Значение/ Диапазон & Единица измерения & Шаг \\ 
\hline
\endfirsthead
\caption*{\hspace{3pt}\emph{Продолжение таблицы \ref{tznp:tbl1}\hfill\vspace{-0.5\baselineskip}}} \\ % сделано по ГОСТ 2.105 п.6.8.7
\hline
\rowcolor{unimain!40}
Параметр (Параметр на ИЧМ) & Условное обозначение на схеме & Значение/ Диапазон & Единица измерения & Шаг \\ 
%\hline
\endhead
%\hline
\endfoot
%\hline
\endlastfoot
%==+t1*ZSPTOC1|SV_LVZSTOC
\multicolumn{5}{|c|}{ ТЗНП 1 ст } \\ \hline 
\centering Ввод функции в работу (Ввод\_функции)  & \centering SGF1 & \centering Не предусмотрено \\ Предусмотрено & \centering -- & \centering \arraybackslash -- \\
\hline
\centering Режим оперативного блокирования ЧС (Реж\_БЧС)  & \centering SGF8 & \centering Не предусмотрено \\ Предусмотрено & \centering -- & \centering \arraybackslash -- \\
\hline
\centering Режим контроля от БНТ (Реж\_БНТ)  & \centering SGF7 & \centering Не предусмотрено \\ Предусмотрено & \centering -- & \centering \arraybackslash -- \\
\hline
\centering Ток срабатывания НП (3I0ср)  & \centering 3I0> & \centering 0,10...30,00 & \centering о.е. & \centering \arraybackslash 0,01 \\
\hline
\centering Вывод направленности при АУ (Ненапр\_АУ)  & \centering SGF4 & \centering Не предусмотрено \\ Предусмотрено & \centering -- & \centering \arraybackslash -- \\
\hline
\centering Режим направленности (Направленность)  & \centering SGF2 & \centering Ненаправленная \\ Прямонаправленная \\ Обратнонаправленная & \centering -- & \centering \arraybackslash -- \\
\hline
\centering Режим ОНМ НП при неисправности ЦН (Реж\_БНН\_ОНМ\_НП)  & \centering SGF3 & \centering Блокировка \\ Деблокировка & \centering -- & \centering \arraybackslash -- \\
\hline
\centering Режим ПО 3U0 при неисправности ЦН (Реж\_БНН\_3U0)  & \centering SGF5 & \centering Блокировка \\ Деблокировка & \centering -- & \centering \arraybackslash -- \\
\hline
\centering Напряжение срабатывания НП (3U0ср)  & \centering 3U0> & \centering 2,0...100,0 & \centering В & \centering \arraybackslash 0,1 \\
\hline
\centering Режим контроля от ПО 3U0 (Реж\_3U0)  & \centering SGF6 & \centering Не предусмотрено \\ Предусмотрено & \centering -- & \centering \arraybackslash -- \\
\hline
\centering Характеристика срабатывания (Характеристика)  & \centering -- & \centering Кривая F \\ Кривая E \\ Кривая D \\ Кривая B \\ Кривая A \\ Кривая C \\ Независимая \\ Кривая RXIDG & \centering -- & \centering \arraybackslash -- \\
\hline
\centering Коэффициент регулировки времени срабатывания зависимой характеристики (К)  & \centering -- & \centering 0,05...15,00 & \centering -- & \centering \arraybackslash 0,01 \\
\hline
\centering Независимая выдержка времени срабатывания (Tср)  & \centering T1 & \centering 0...100000 & \centering мс & \centering \arraybackslash 10 \\
\hline
\multicolumn{5}{|c|}{ ТЗНП 2 ст } \\ \hline 
\centering Ввод функции в работу (Ввод\_функции)  & \centering SGF1 & \centering Не предусмотрено \\ Предусмотрено & \centering -- & \centering \arraybackslash -- \\
\hline
\centering Режим оперативного блокирования ЧС (Реж\_БЧС)  & \centering SGF8 & \centering Не предусмотрено \\ Предусмотрено & \centering -- & \centering \arraybackslash -- \\
\hline
\centering Режим контроля от БНТ (Реж\_БНТ)  & \centering SGF7 & \centering Не предусмотрено \\ Предусмотрено & \centering -- & \centering \arraybackslash -- \\
\hline
\centering Ток срабатывания НП (3I0ср)  & \centering 3I0> & \centering 0,10...30,00 & \centering о.е. & \centering \arraybackslash 0,01 \\
\hline
\centering Вывод направленности при АУ (Ненапр\_АУ)  & \centering SGF4 & \centering Не предусмотрено \\ Предусмотрено & \centering -- & \centering \arraybackslash -- \\
\hline
\centering Режим направленности (Направленность)  & \centering SGF2 & \centering Ненаправленная \\ Прямонаправленная \\ Обратнонаправленная & \centering -- & \centering \arraybackslash -- \\
\hline
\centering Режим ОНМ НП при неисправности ЦН (Реж\_БНН\_ОНМ\_НП)  & \centering SGF3 & \centering Блокировка \\ Деблокировка & \centering -- & \centering \arraybackslash -- \\
\hline
\centering Режим ПО 3U0 при неисправности ЦН (Реж\_БНН\_3U0)  & \centering SGF5 & \centering Блокировка \\ Деблокировка & \centering -- & \centering \arraybackslash -- \\
\hline
\centering Напряжение срабатывания НП (3U0ср)  & \centering 3U0> & \centering 2,0...100,0 & \centering В & \centering \arraybackslash 0,1 \\
\hline
\centering Режим контроля от ПО 3U0 (Реж\_3U0)  & \centering SGF6 & \centering Не предусмотрено \\ Предусмотрено & \centering -- & \centering \arraybackslash -- \\
\hline
\centering Характеристика срабатывания (Характеристика)  & \centering -- & \centering Кривая F \\ Кривая E \\ Кривая D \\ Кривая B \\ Кривая A \\ Кривая C \\ Независимая \\ Кривая RXIDG & \centering -- & \centering \arraybackslash -- \\
\hline
\centering Коэффициент регулировки времени срабатывания зависимой характеристики (К)  & \centering -- & \centering 0,05...15,00 & \centering -- & \centering \arraybackslash 0,01 \\
\hline
\centering Независимая выдержка времени срабатывания (Tср)  & \centering T1 & \centering 0...100000 & \centering мс & \centering \arraybackslash 10 \\
\hline
\multicolumn{5}{|c|}{ ТЗНП 3 ст } \\ \hline 
\centering Ввод функции в работу (Ввод\_функции)  & \centering SGF1 & \centering Не предусмотрено \\ Предусмотрено & \centering -- & \centering \arraybackslash -- \\
\hline
\centering Режим оперативного блокирования ЧС (Реж\_БЧС)  & \centering SGF8 & \centering Не предусмотрено \\ Предусмотрено & \centering -- & \centering \arraybackslash -- \\
\hline
\centering Режим контроля от БНТ (Реж\_БНТ)  & \centering SGF7 & \centering Не предусмотрено \\ Предусмотрено & \centering -- & \centering \arraybackslash -- \\
\hline
\centering Ток срабатывания НП (3I0ср)  & \centering 3I0> & \centering 0,10...30,00 & \centering о.е. & \centering \arraybackslash 0,01 \\
\hline
\centering Вывод направленности при АУ (Ненапр\_АУ)  & \centering SGF4 & \centering Не предусмотрено \\ Предусмотрено & \centering -- & \centering \arraybackslash -- \\
\hline
\centering Режим направленности (Направленность)  & \centering SGF2 & \centering Ненаправленная \\ Прямонаправленная \\ Обратнонаправленная & \centering -- & \centering \arraybackslash -- \\
\hline
\centering Режим ОНМ НП при неисправности ЦН (Реж\_БНН\_ОНМ\_НП)  & \centering SGF3 & \centering Блокировка \\ Деблокировка & \centering -- & \centering \arraybackslash -- \\
\hline
\centering Режим ПО 3U0 при неисправности ЦН (Реж\_БНН\_3U0)  & \centering SGF5 & \centering Блокировка \\ Деблокировка & \centering -- & \centering \arraybackslash -- \\
\hline
\centering Напряжение срабатывания НП (3U0ср)  & \centering 3U0> & \centering 2,0...100,0 & \centering В & \centering \arraybackslash 0,1 \\
\hline
\centering Режим контроля от ПО 3U0 (Реж\_3U0)  & \centering SGF6 & \centering Не предусмотрено \\ Предусмотрено & \centering -- & \centering \arraybackslash -- \\
\hline
\centering Характеристика срабатывания (Характеристика)  & \centering -- & \centering Кривая F \\ Кривая E \\ Кривая D \\ Кривая B \\ Кривая A \\ Кривая C \\ Независимая \\ Кривая RXIDG & \centering -- & \centering \arraybackslash -- \\
\hline
\centering Коэффициент регулировки времени срабатывания зависимой характеристики (К)  & \centering -- & \centering 0,05...15,00 & \centering -- & \centering \arraybackslash 0,01 \\
\hline
\centering Независимая выдержка времени срабатывания (Tср)  & \centering T1 & \centering 0...100000 & \centering мс & \centering \arraybackslash 10 \\
\hline
\multicolumn{5}{|c|}{ ТЗНП 4 ст } \\ \hline 
\centering Ввод функции в работу (Ввод\_функции)  & \centering SGF1 & \centering Не предусмотрено \\ Предусмотрено & \centering -- & \centering \arraybackslash -- \\
\hline
\centering Режим оперативного блокирования ЧС (Реж\_БЧС)  & \centering SGF8 & \centering Не предусмотрено \\ Предусмотрено & \centering -- & \centering \arraybackslash -- \\
\hline
\centering Режим контроля от БНТ (Реж\_БНТ)  & \centering SGF7 & \centering Не предусмотрено \\ Предусмотрено & \centering -- & \centering \arraybackslash -- \\
\hline
\centering Ток срабатывания НП (3I0ср)  & \centering 3I0> & \centering 0,10...30,00 & \centering о.е. & \centering \arraybackslash 0,01 \\
\hline
\centering Вывод направленности при АУ (Ненапр\_АУ)  & \centering SGF4 & \centering Не предусмотрено \\ Предусмотрено & \centering -- & \centering \arraybackslash -- \\
\hline
\centering Режим направленности (Направленность)  & \centering SGF2 & \centering Ненаправленная \\ Прямонаправленная \\ Обратнонаправленная & \centering -- & \centering \arraybackslash -- \\
\hline
\centering Режим ОНМ НП при неисправности ЦН (Реж\_БНН\_ОНМ\_НП)  & \centering SGF3 & \centering Блокировка \\ Деблокировка & \centering -- & \centering \arraybackslash -- \\
\hline
\centering Режим ПО 3U0 при неисправности ЦН (Реж\_БНН\_3U0)  & \centering SGF5 & \centering Блокировка \\ Деблокировка & \centering -- & \centering \arraybackslash -- \\
\hline
\centering Напряжение срабатывания НП (3U0ср)  & \centering 3U0> & \centering 2,0...100,0 & \centering В & \centering \arraybackslash 0,1 \\
\hline
\centering Режим контроля от ПО 3U0 (Реж\_3U0)  & \centering SGF6 & \centering Не предусмотрено \\ Предусмотрено & \centering -- & \centering \arraybackslash -- \\
\hline
\centering Характеристика срабатывания (Характеристика)  & \centering -- & \centering Кривая F \\ Кривая E \\ Кривая D \\ Кривая B \\ Кривая A \\ Кривая C \\ Независимая \\ Кривая RXIDG & \centering -- & \centering \arraybackslash -- \\
\hline
\centering Коэффициент регулировки времени срабатывания зависимой характеристики (К)  & \centering -- & \centering 0,05...15,00 & \centering -- & \centering \arraybackslash 0,01 \\
\hline
\centering Независимая выдержка времени срабатывания (Tср)  & \centering T1 & \centering 0...100000 & \centering мс & \centering \arraybackslash 10 \\
\hline
\multicolumn{5}{|c|}{ АУ ТЗНП } \\ \hline 
\centering Выбор ускоряемой ступени (Выбор\_ст)  & \centering SGF1 & \centering 1 ступень \\ 2 ступень \\ 3 ступень \\ 4 ступень \\ 1 или 3 ступень \\ 2 или 4 ступень & \centering -- & \centering \arraybackslash -- \\
\hline
\multicolumn{5}{|c|}{ ОНМ НП } \\ \hline 
\centering Угол максимальной чувствительности (Фмч)  & \centering φмч & \centering -179,0...180,0 & \centering град. & \centering \arraybackslash 0,1 \\
\hline
\multicolumn{5}{|c|}{ БНТ НП } \\ \hline 
\centering Отношение тока второй гармоники к току основной гармоники (Ih2/Ih1)  & \centering I2г/I1г> & \centering 5...100 & \centering \% & \centering \arraybackslash 1 \\
\hline
\centering Ток блокировки НП (3I0ср)  & \centering 3I0> & \centering 0,1...15,0 & \centering о.е. & \centering \arraybackslash 0,1 \\
\hline
\multicolumn{5}{|c|}{ ОУ ТЗНП } \\ \hline 
\centering Выбор оперативно ускоряемой ступени (Выбор\_ст\_ОУ)  & \centering SGF1 & \centering 1 ступень \\ 2 ступень \\ 3 ступень \\ 4 ступень \\ 1 или 3 ступень \\ 2 или 4 ступень & \centering -- & \centering \arraybackslash -- \\
\hline
\centering Выдержка времени срабатывания оперативного ускорения (T\_ОУ)  & \centering T1 & \centering 0...3000 & \centering мс & \centering \arraybackslash 10 \\
\hline
%===t1
\end{longtable}
\normalsize



%%%%%%%%%%%%%%%%%%%%%%%%%%%%%%%%%%%%%%%%%%%%%%%%%%%%%%%%%% УРОВ %%%%%%%%%%%%%%%%%%%%%%%%%%%%%%%%%%%%%%%%%%%%%%%%%%%%%%%%%%%%%%%%%%
\par
В таблице \ref{urov35:tbl1} приведены параметры, необходимые для настройки УРОВ.

\footnotesize
\begin{longtable}{|>{\centering\arraybackslash}m{5.3cm}|>{\centering\arraybackslash}m{3.3cm}|>{\centering\arraybackslash}m{4.2cm}|>{\centering\arraybackslash}m{1.8cm}|>{\centering\arraybackslash}m{1cm}|}
\caption{Параметры для настройки УРОВ\hfill\vspace{-0.5\baselineskip}}\label{urov35:tbl1}\\ 
\hline
\rowcolor{gray!20}
Параметр (Параметр на ИЧМ) & Условное обозначение на схеме & Значение/ Диапазон & Единица измерения & Шаг \\ 
\hline
\endfirsthead
\caption*{\hspace{3pt}\emph{Продолжение таблицы \ref{urov35:tbl1}\hfill\vspace{-0.5\baselineskip}}} \\ % сделано по ГОСТ 2.105 п.6.8.7
\hline
\rowcolor{gray!20}
Параметр (Параметр на ИЧМ) & Условное обозначение на схеме & Значение/ Диапазон & Единица измерения & Шаг \\ 
%\hline
\endhead
%\hline
\endfoot
%\hline
\endlastfoot
%==+t1*GENRBRF1|SV_TPBRF
\centering Ввод функции в работу (Ввод\_функции)  & \centering SGF1 & \centering Не предусмотрено \\ Предусмотрено & \centering -- & \centering \arraybackslash -- \\
\hline
\centering Ток срабатывания (Iср)  & \centering I> & \centering 0,05...0,50 & \centering о.е. & \centering \arraybackslash 0,01 \\
\hline
\centering Ускорение при блокировке отключения В (Уск\_при\_блк\_откл\_В)  & \centering SGF2 & \centering Не предусмотрено \\ Предусмотрено & \centering -- & \centering \arraybackslash -- \\
\hline
\centering УРОВ с подхватом по току (Подхват\_по\_току)  & \centering SGF3 & \centering Не предусмотрено \\ Предусмотрено & \centering -- & \centering \arraybackslash -- \\
\hline
\centering Контроль по току при действии "на себя" (Контр\_тока\_на\_себя)  & \centering SGF6 & \centering Не предусмотрено \\ Предусмотрено по внутр. ПО \\ Предусмотрено по внеш. ПО & \centering -- & \centering \arraybackslash -- \\
\hline
\centering Контроль ЭМО (Контроль\_ЭМО)  & \centering SGF4 & \centering Не предусмотрено \\ Предусмотрено по ЭМО1 \\ Предусмотрено по ЭМО1 и ЭМО2 & \centering -- & \centering \arraybackslash -- \\
\hline
\centering Действие внешнего УРОВ на вышестоящий выключатель (Действ\_на\_выш\_выкл)  & \centering SGF5 & \centering Не предусмотрено \\ Предусмотрено & \centering -- & \centering \arraybackslash -- \\
\hline
\centering Выдержка времени срабатывания (Tср)  & \centering T1 & \centering 0...1000 & \centering мс & \centering \arraybackslash 10 \\
\hline
%===t1
\end{longtable}
\normalsize


%%%%%%%%%%%%%%%%%%%%%%%%%%%%%%%%%%%%%%%%%%%%%%%%%%%%%%%%%% БНН 1сш %%%%%%%%%%%%%%%%%%%%%%%%%%%%%%%%%%%%%%%%%%%%%%%%%%%%%%%%%%%%%%%%%%
\par
В таблице \ref{bnnsh1:tbl1} приведены параметры, необходимые для настройки БНН 1сш.

\footnotesize
\begin{longtable}{|>{\centering\arraybackslash}m{5.3cm}|>{\centering\arraybackslash}m{3.3cm}|>{\centering\arraybackslash}m{4.2cm}|>{\centering\arraybackslash}m{1.8cm}|>{\centering\arraybackslash}m{1cm}|}
\caption{Параметры для настройки БНН 1сш\hfill\vspace{-0.5\baselineskip}}\label{bnnsh1:tbl1}\\ 
\hline
\rowcolor{unimain!40}
Параметр (Параметр на ИЧМ) & Условное обозначение на схеме & Значение/ Диапазон & Единица измерения & Шаг \\ 
\hline
\endfirsthead
\caption*{\hspace{3pt}\emph{Продолжение таблицы \ref{bnnsh1:tbl1}\hfill\vspace{-0.5\baselineskip}}} \\ % сделано по ГОСТ 2.105 п.6.8.7
\hline
\rowcolor{unimain!40}
Параметр (Параметр на ИЧМ) & Условное обозначение на схеме & Значение/ Диапазон & Единица измерения & Шаг \\ 
%\hline
\endhead
%\hline
\endfoot
%\hline
\endlastfoot
%==+t1*RVTR1|SV_LVVTR1
\centering Ввод функции в работу (Ввод\_функции)  & \centering SGF1 & \centering Не предусмотрено \\ Предусмотрено & \centering -- & \centering \arraybackslash -- \\
\hline
\centering Уровень небаланса по напряжению (dU0ср)  & \centering ∑Uф--3U0> & \centering 1,0...150,0 & \centering В & \centering \arraybackslash 0,1 \\
\hline
\centering Ток срабатывания ОП (I2ср)  & \centering I2< & \centering 0,04...4,00 & \centering о.е. & \centering \arraybackslash 0,01 \\
\hline
\centering Напряжение срабатывания ОП (U2ср)  & \centering U2> & \centering 2,0...60,0 & \centering В & \centering \arraybackslash 0,1 \\
\hline
\centering Напряжение срабатывания (Uср)  & \centering U< & \centering 2,0...100,0 & \centering В & \centering \arraybackslash 0,1 \\
\hline
\centering Ток срабатывания (Iср)  & \centering I> & \centering 0,04...4,00 & \centering о.е. & \centering \arraybackslash 0,01 \\
\hline
\centering Напряжение срабатывания НП третьей гармоники (3U0ср\_3гарм)  & \centering 3U03f< & \centering 0,05...10,00 & \centering В & \centering \arraybackslash 0,01 \\
\hline
\centering Ввод в работу контроля небаланса (Контроль\_dU0)  & \centering SGF2 & \centering Не предусмотрено \\ Предусмотрено & \centering -- & \centering \arraybackslash -- \\
\hline
\centering Ввод в работу контроля напряжения ОП (Контроль\_U2)  & \centering SGF3 & \centering Не предусмотрено \\ Предусмотрено & \centering -- & \centering \arraybackslash -- \\
\hline
\centering Ввод в работу контроля минимального напряжения (Контроль\_Uмин)  & \centering SGF4 & \centering Не предусмотрено \\ Предусмотрено & \centering -- & \centering \arraybackslash -- \\
\hline
\centering Ввод в работу контроля 3 гармоники напряжения НП (Контроль\_3гарм\_3U0)  & \centering SGF5 & \centering Не предусмотрено \\ Предусмотрено & \centering -- & \centering \arraybackslash -- \\
\hline
\centering Выдержка времени срабатывания сигнализации (Tсигн)  & \centering T1 & \centering 0...100000 & \centering мс & \centering \arraybackslash 10 \\
\hline
%===t1
\end{longtable}
\normalsize


%%%%%%%%%%%%%%%%%%%%%%%%%%%%%%%%%%%%%%%%%%%%%%%%%%%%%%%%%% БНН 2сш %%%%%%%%%%%%%%%%%%%%%%%%%%%%%%%%%%%%%%%%%%%%%%%%%%%%%%%%%%%%%%%%%%
\par
В таблице \ref{bnnsh2:tbl1} приведены параметры, необходимые для настройки БНН 2сш.

\footnotesize
\begin{longtable}{|>{\centering\arraybackslash}m{5.3cm}|>{\centering\arraybackslash}m{3.3cm}|>{\centering\arraybackslash}m{4.2cm}|>{\centering\arraybackslash}m{1.8cm}|>{\centering\arraybackslash}m{1cm}|}
\caption{Параметры для настройки БНН 2сш\hfill\vspace{-0.5\baselineskip}}\label{bnnsh2:tbl1}\\ 
\hline
\rowcolor{unimain!40}
Параметр (Параметр на ИЧМ) & Условное обозначение на схеме & Значение/ Диапазон & Единица измерения & Шаг \\ 
\hline
\endfirsthead
\caption*{\hspace{3pt}\emph{Продолжение таблицы \ref{bnnsh2:tbl1}\hfill\vspace{-0.5\baselineskip}}} \\ % сделано по ГОСТ 2.105 п.6.8.7
\hline
\rowcolor{unimain!40}
Параметр (Параметр на ИЧМ) & Условное обозначение на схеме & Значение/ Диапазон & Единица измерения & Шаг \\ 
%\hline
\endhead
%\hline
\endfoot
%\hline
\endlastfoot
%==+t1*RVTR1|SV_LVVTR2
\centering Ввод функции в работу (Ввод\_функции)  & \centering SGF1 & \centering Не предусмотрено \\ Предусмотрено & \centering -- & \centering \arraybackslash -- \\
\hline
\centering Ток срабатывания ОП (I2ср)  & \centering I2< & \centering 0,04...4,00 & \centering о.е. & \centering \arraybackslash 0,01 \\
\hline
\centering Напряжение срабатывания ОП (U2ср)  & \centering U2> & \centering 2,0...60,0 & \centering В & \centering \arraybackslash 0,1 \\
\hline
\centering Напряжение срабатывания (Uср)  & \centering U< & \centering 2,0...100,0 & \centering В & \centering \arraybackslash 0,1 \\
\hline
\centering Ток срабатывания (Iср)  & \centering I> & \centering 0,04...4,00 & \centering о.е. & \centering \arraybackslash 0,01 \\
\hline
\centering Ввод в работу контроля напряжения ОП (Контроль\_U2)  & \centering SGF2 & \centering Не предусмотрено \\ Предусмотрено & \centering -- & \centering \arraybackslash -- \\
\hline
\centering Ввод в работу контроля минимального напряжения (Контроль\_Uмин)  & \centering SGF3 & \centering Не предусмотрено \\ Предусмотрено & \centering -- & \centering \arraybackslash -- \\
\hline
\centering Выдержка времени срабатывания сигнализации (Tсигн)  & \centering T1 & \centering 0...100000 & \centering мс & \centering \arraybackslash 10 \\
\hline
%===t1
\end{longtable}
\normalsize


%%%%%%%%%%%%%%%%%%%%%%%%%%%%%%%%%%%%%%%%%%%%%%%%%%%%%%%%%% КНН 1сш %%%%%%%%%%%%%%%%%%%%%%%%%%%%%%%%%%%%%%%%%%%%%%%%%%%%%%%%%%%%%%%%%%
\par
В таблице \ref{knnsh:tbl1} приведены параметры, необходимые для настройки КНН 1сш.

\footnotesize
\begin{longtable}{|>{\centering\arraybackslash}m{5.3cm}|>{\centering\arraybackslash}m{3.3cm}|>{\centering\arraybackslash}m{4.2cm}|>{\centering\arraybackslash}m{1.8cm}|>{\centering\arraybackslash}m{1cm}|}
\caption{Параметры для настройки функции <<КНН 1(2)сш>>\hfill\vspace{-0.5\baselineskip}}\label{knnsh:tbl1}\\ 
\hline
\rowcolor{unimain!40}
Параметр (Параметр на ИЧМ) & Условное обозначение на схеме & Значение/ Диапазон & Единица измерения & Шаг \\ 
\hline
\endfirsthead
\caption*{\hspace{3pt}\emph{Продолжение таблицы \ref{knnsh:tbl1}\hfill\vspace{-0.5\baselineskip}}} \\ % сделано по ГОСТ 2.105 п.6.8.7
\hline
\rowcolor{unimain!40}
Параметр (Параметр на ИЧМ) & Условное обозначение на схеме & Значение/ Диапазон & Единица измерения & Шаг \\ 
%\hline
\endhead
%\hline
\endfoot
%\hline
\endlastfoot
%==+t1*PTOV1|SV_VOLPRES1
\centering Ввод функции в работу (Ввод\_функции)  & \centering SGF1 & \centering Не предусмотрено \\ Предусмотрено & \centering -- & \centering \arraybackslash -- \\
\hline
\centering Напряжение срабатывания (Uср)  & \centering U> & \centering 50,0...150,0 & \centering В & \centering \arraybackslash 0,1 \\
\hline
\centering Напряжение срабатывания ОП (U2ср)  & \centering U2> & \centering 2,0...50,0 & \centering В & \centering \arraybackslash 0,1 \\
\hline
\centering Режим пуска (Реж\_пуска)  & \centering SGF2 & \centering Измеренный \\ Внешний & \centering -- & \centering \arraybackslash -- \\
\hline
%===t1
\end{longtable}
\normalsize


%%%%%%%%%%%%%%%%%%%%%%%%%%%%%%%%%%%%%%%%%%%%%%%%%%%%%%%%%% КНН 2сш %%%%%%%%%%%%%%%%%%%%%%%%%%%%%%%%%%%%%%%%%%%%%%%%%%%%%%%%%%%%%%%%%%
\par
В таблице \ref{knnsh1:tbl1} приведены параметры, необходимые для настройки КНН 2сш.

\footnotesize
\begin{longtable}{|>{\centering\arraybackslash}m{5.3cm}|>{\centering\arraybackslash}m{3.3cm}|>{\centering\arraybackslash}m{4.2cm}|>{\centering\arraybackslash}m{1.8cm}|>{\centering\arraybackslash}m{1cm}|}
\caption{Параметры для настройки функции <<КНН 1(2)сш>>\hfill\vspace{-0.5\baselineskip}}\label{knnsh1:tbl1}\\ 
\hline
\rowcolor{unimain!40}
Параметр (Параметр на ИЧМ) & Условное обозначение на схеме & Значение/ Диапазон & Единица измерения & Шаг \\ 
\hline
\endfirsthead
\caption*{\hspace{3pt}\emph{Продолжение таблицы \ref{knnsh1:tbl1}\hfill\vspace{-0.5\baselineskip}}} \\ % сделано по ГОСТ 2.105 п.6.8.7
\hline
\rowcolor{unimain!40}
Параметр (Параметр на ИЧМ) & Условное обозначение на схеме & Значение/ Диапазон & Единица измерения & Шаг \\ 
%\hline
\endhead
%\hline
\endfoot
%\hline
\endlastfoot
%==+t1*PTOV1|SV_VOLPRES2
\centering Ввод функции в работу (Ввод\_функции)  & \centering SGF1 & \centering Не предусмотрено \\ Предусмотрено & \centering -- & \centering \arraybackslash -- \\
\hline
\centering Напряжение срабатывания (Uср)  & \centering U> & \centering 50,0...150,0 & \centering В & \centering \arraybackslash 0,1 \\
\hline
\centering Напряжение срабатывания ОП (U2ср)  & \centering U2> & \centering 2,0...50,0 & \centering В & \centering \arraybackslash 0,1 \\
\hline
\centering Режим пуска (Реж\_пуска)  & \centering SGF2 & \centering Измеренный \\ Внешний & \centering -- & \centering \arraybackslash -- \\
\hline
%===t1
\end{longtable}
\normalsize



%%%%%%%%%%%%%%%%%%%%%%%%%%%%%%%%%%%%%%%%%%%%%%%%%%%%%%%%%% КОН 1сш %%%%%%%%%%%%%%%%%%%%%%%%%%%%%%%%%%%%%%%%%%%%%%%%%%%%%%%%%%%%%%%%%%
\par
В таблице \ref{konsh:tbl1} приведены параметры, необходимые для настройки КОН 1сш.

\footnotesize
\begin{longtable}{|>{\centering\arraybackslash}m{5.3cm}|>{\centering\arraybackslash}m{3.3cm}|>{\centering\arraybackslash}m{4.2cm}|>{\centering\arraybackslash}m{1.8cm}|>{\centering\arraybackslash}m{1cm}|}
\caption{Параметры для настройки функции «КОН 1(2)сш»\hfill\vspace{-0.5\baselineskip}}\label{konsh:tbl1}\\ 
\hline
\rowcolor{unimain!40}
Параметр (Параметр на ИЧМ) & Условное обозначение на схеме & Значение/ Диапазон & Единица измерения & Шаг \\ 
\hline
\endfirsthead
\caption*{\hspace{3pt}\emph{Продолжение таблицы \ref{konsh:tbl1}\hfill\vspace{-0.5\baselineskip}}} \\ % сделано по ГОСТ 2.105 п.6.8.7
\hline
\rowcolor{unimain!40}
Параметр (Параметр на ИЧМ) & Условное обозначение на схеме & Значение/ Диапазон & Единица измерения & Шаг \\ 
%\hline
\endhead
%\hline
\endfoot
%\hline
\endlastfoot
%==+t1*PTUV1|SV_VOLABS1
\centering Ввод функции в работу (Ввод\_функции)  & \centering SGF1 & \centering Не предусмотрено \\ Предусмотрено & \centering -- & \centering \arraybackslash -- \\
\hline
\centering Напряжение срабатывания (Uср)  & \centering U< & \centering 2,0...100,0 & \centering В & \centering \arraybackslash 0,1 \\
\hline
\centering Напряжение срабатывания ОП (U2ср)  & \centering U2> & \centering 2,0...50,0 & \centering В & \centering \arraybackslash 0,1 \\
\hline
\centering Режим пуска (Реж\_пуска)  & \centering SGF2 & \centering Измеренный \\ Внешний & \centering -- & \centering \arraybackslash -- \\
\hline
%===t1
\end{longtable}
\normalsize


%%%%%%%%%%%%%%%%%%%%%%%%%%%%%%%%%%%%%%%%%%%%%%%%%%%%%%%%%% КОН 2сш %%%%%%%%%%%%%%%%%%%%%%%%%%%%%%%%%%%%%%%%%%%%%%%%%%%%%%%%%%%%%%%%%%
\par
В таблице \ref{konsh1:tbl1} приведены параметры, необходимые для настройки КОН 2сш.

\footnotesize
\begin{longtable}{|>{\centering\arraybackslash}m{5.3cm}|>{\centering\arraybackslash}m{3.3cm}|>{\centering\arraybackslash}m{4.2cm}|>{\centering\arraybackslash}m{1.8cm}|>{\centering\arraybackslash}m{1cm}|}
\caption{Параметры для настройки функции «КОН 1(2)сш»\hfill\vspace{-0.5\baselineskip}}\label{konsh1:tbl1}\\ 
\hline
\rowcolor{unimain!40}
Параметр (Параметр на ИЧМ) & Условное обозначение на схеме & Значение/ Диапазон & Единица измерения & Шаг \\ 
\hline
\endfirsthead
\caption*{\hspace{3pt}\emph{Продолжение таблицы \ref{konsh1:tbl1}\hfill\vspace{-0.5\baselineskip}}} \\ % сделано по ГОСТ 2.105 п.6.8.7
\hline
\rowcolor{unimain!40}
Параметр (Параметр на ИЧМ) & Условное обозначение на схеме & Значение/ Диапазон & Единица измерения & Шаг \\ 
%\hline
\endhead
%\hline
\endfoot
%\hline
\endlastfoot
%==+t1*PTUV1|SV_VOLABS2
\centering Ввод функции в работу (Ввод\_функции)  & \centering SGF1 & \centering Не предусмотрено \\ Предусмотрено & \centering -- & \centering \arraybackslash -- \\
\hline
\centering Напряжение срабатывания (Uср)  & \centering U< & \centering 2,0...100,0 & \centering В & \centering \arraybackslash 0,1 \\
\hline
\centering Напряжение срабатывания ОП (U2ср)  & \centering U2> & \centering 2,0...50,0 & \centering В & \centering \arraybackslash 0,1 \\
\hline
\centering Режим пуска (Реж\_пуска)  & \centering SGF2 & \centering Измеренный \\ Внешний & \centering -- & \centering \arraybackslash -- \\
\hline
%===t1
\end{longtable}
\normalsize


%%%%%%%%%%%%%%%%%%%%%%%%%%%%%%%%%%%%%%%%%%%%%%%%%%%%%%%%%% КС %%%%%%%%%%%%%%%%%%%%%%%%%%%%%%%%%%%%%%%%%%%%%%%%%%%%%%%%%%%%%%%%%%
\par
В таблице \ref{ks:tbl1} приведены параметры, необходимые для настройки КС.

\footnotesize
\begin{longtable}{|>{\centering\arraybackslash}m{5.3cm}|>{\centering\arraybackslash}m{3.3cm}|>{\centering\arraybackslash}m{4.2cm}|>{\centering\arraybackslash}m{1.8cm}|>{\centering\arraybackslash}m{1cm}|}
\caption{Параметры для настройки КС\hfill\vspace{-0.5\baselineskip}}\label{ks:tbl1}\\ 
\hline
\rowcolor{unimain!40}
Параметр (Параметр на ИЧМ) & Условное обозначение на схеме & Значение/ Диапазон & Единица измерения & Шаг \\ 
\hline
\endfirsthead
\caption*{\hspace{3pt}\emph{Продолжение таблицы \ref{ks:tbl1}\hfill\vspace{-0.5\baselineskip}}} \\ % сделано по ГОСТ 2.105 п.6.8.7
\hline
\rowcolor{unimain!40}
Параметр (Параметр на ИЧМ) & Условное обозначение на схеме & Значение/ Диапазон & Единица измерения & Шаг \\ 
%\hline
\endhead
%\hline
\endfoot
%\hline
\endlastfoot
%==+t1*LVRSYN1|SV_LVSYN
\centering Ввод функции в работу (Ввод\_функции)  & \centering SGF1 & \centering Не предусмотрено \\ Предусмотрено & \centering -- & \centering \arraybackslash -- \\
\hline
\centering Синхронизируемое напряжение (Синхрон\_напряж)  & \centering SGF2 & \centering UAB \\ UBC \\ UCA & \centering -- & \centering \arraybackslash -- \\
\hline
\centering Допустимая разность модуля векторов напряжений (dUср)  & \centering dU< & \centering 5,0...50,0 & \centering В & \centering \arraybackslash 0,1 \\
\hline
\centering Допустимая разность частот для ожидания синхронизма (dFср)  & \centering dFос< & \centering 0,05...0,50 & \centering Гц & \centering \arraybackslash 0,01 \\
\hline
\centering Разность частот срабатывания и возврата ИО частоты ожидания синхронизма (dFвоз)  & \centering -- & \centering 0,04...0,10 & \centering Гц & \centering \arraybackslash 0,01 \\
\hline
\centering Максимальная разность частот для улавливания синхронизма (dFмакс\_ср)  & \centering dFмакс< & \centering 0,05...2,00 & \centering Гц & \centering \arraybackslash 0,01 \\
\hline
\centering Разность частот срабатывания и возврата ИО частоты улавливания синхронизма (dFмакс\_воз)  & \centering -- & \centering 0,04...0,10 & \centering Гц & \centering \arraybackslash 0,01 \\
\hline
\centering Максимальная скорость изменения частоты (dFdtср)  & \centering l dF/dt l> & \centering 0,1...15,0 & \centering Гц/с & \centering \arraybackslash 0,1 \\
\hline
\centering Коэффициент возврата по скорости изменения частоты (Kвоз)  & \centering -- & \centering 0,20...0,99 & \centering -- & \centering \arraybackslash 0,01 \\
\hline
\centering Допустимая разность фаз векторов напряжений (dФср)  & \centering dФос< & \centering 10,0...90,0 & \centering град. & \centering \arraybackslash 0,1 \\
\hline
\centering Угол сдвига (Ф\_сдвига)  & \centering Ф\_сдвига & \centering 0,0...359,9 & \centering град. & \centering \arraybackslash 0,1 \\
\hline
\centering Улавливание синхронизма (Улавл\_синхр)  & \centering SGF3 & \centering Не предусмотрено \\ Предусмотрено & \centering -- & \centering \arraybackslash -- \\
\hline
\centering Максимально допустимый угол включения выключателя (dФус)  & \centering dФус< & \centering 10,0...90,0 & \centering град. & \centering \arraybackslash 0,1 \\
\hline
\centering Время включения выключателя (Tопер)  & \centering Tопер & \centering 10...2000 & \centering мс & \centering \arraybackslash 10 \\
\hline
%===t1
\end{longtable}
\normalsize



%%%%%%%%%%%%%%%%%%%%%%%%%%%%%%%%%%%%%%%%%%%%%%%%%%%%%%%%%% АВР %%%%%%%%%%%%%%%%%%%%%%%%%%%%%%%%%%%%%%%%%%%%%%%%%%%%%%%%%%%%%%%%%%
\par
В таблице \ref{avr:tbl1} приведены параметры, необходимые для настройки АВР.

\footnotesize
\begin{longtable}{|>{\centering\arraybackslash}m{5.3cm}|>{\centering\arraybackslash}m{3.3cm}|>{\centering\arraybackslash}m{4.2cm}|>{\centering\arraybackslash}m{1.8cm}|>{\centering\arraybackslash}m{1cm}|}
\caption{Параметры для настройки АВР\hfill\vspace{-0.5\baselineskip}}\label{avr:tbl1}\\ 
\hline
\rowcolor{unimain!40}
Параметр (Параметр на ИЧМ) & Условное обозначение на схеме & Значение/ Диапазон & Единица измерения & Шаг \\ 
\hline
\endfirsthead
\caption*{\hspace{3pt}\emph{Продолжение таблицы \ref{avr:tbl1}\hfill\vspace{-0.5\baselineskip}}} \\ % сделано по ГОСТ 2.105 п.6.8.7
\hline
\rowcolor{unimain!40}
Параметр (Параметр на ИЧМ) & Условное обозначение на схеме & Значение/ Диапазон & Единица измерения & Шаг \\ 
%\hline
\endhead
%\hline
\endfoot
%\hline
\endlastfoot
%==+t1*ABTS1|LVBTS
\multicolumn{5}{|c|}{ ЗАВР } \\ \hline 
\centering Режим запрета от СЗЗ (ЗапАВР\_СЗЗ)  & \centering SGF1 & \centering Не предусмотрено \\ Предусмотрено & \centering -- & \centering \arraybackslash -- \\
\hline
\multicolumn{5}{|c|}{ АВР } \\ \hline 
\centering Ввод функции в работу (Ввод\_функции)  & \centering SGF1 & \centering Не предусмотрено \\ Предусмотрено & \centering -- & \centering \arraybackslash -- \\
\hline
\centering Режим работы АВР (Реж\_АВР)  & \centering SGF2 & \centering Внутренний \\ Внешний & \centering -- & \centering \arraybackslash -- \\
\hline
\centering Режим контроля отсутствия напряжения на секции (Реж\_КОН)  & \centering SGF3 & \centering Без контроля \\ С контролем & \centering -- & \centering \arraybackslash -- \\
\hline
\centering Выдержка времени готовности (Тгот)  & \centering T2 & \centering 100...100000 & \centering мс & \centering \arraybackslash 10 \\
\hline
\centering Выдержка времени срабатывания (Tср)  & \centering T1 & \centering 100...10000 & \centering мс & \centering \arraybackslash 10 \\
\hline
%===t1
\end{longtable}
\normalsize


%%%%%%%%%%%%%%%%%%%%%%%%%%%%%%%%%%%%%%%%%%%%%%%%%%%%%%%%%% ВНР %%%%%%%%%%%%%%%%%%%%%%%%%%%%%%%%%%%%%%%%%%%%%%%%%%%%%%%%%%%%%%%%%%
\par
В таблице \ref{vnr:tbl1} приведены параметры, необходимые для настройки ВНР.

\footnotesize
\begin{longtable}{|>{\centering\arraybackslash}m{5.3cm}|>{\centering\arraybackslash}m{3.3cm}|>{\centering\arraybackslash}m{4.2cm}|>{\centering\arraybackslash}m{1.8cm}|>{\centering\arraybackslash}m{1cm}|}
\caption{Параметры для настройки ВНР\hfill\vspace{-0.5\baselineskip}}\label{vnr:tbl1}\\ 
\hline
\rowcolor{unimain!40}
Параметр (Параметр на ИЧМ) & Условное обозначение на схеме & Значение/ Диапазон & Единица измерения & Шаг \\ 
\hline
\endfirsthead
\caption*{\hspace{3pt}\emph{Продолжение таблицы \ref{vnr:tbl1}\hfill\vspace{-0.5\baselineskip}}} \\ % сделано по ГОСТ 2.105 п.6.8.7
\hline
\rowcolor{unimain!40}
Параметр (Параметр на ИЧМ) & Условное обозначение на схеме & Значение/ Диапазон & Единица измерения & Шаг \\ 
%\hline
\endhead
%\hline
\endfoot
%\hline
\endlastfoot
%==+t1*ANSR1|LVNSR
\centering Ввод функции в работу (Ввод\_функции)  & \centering SGF1 & \centering Не предусмотрено \\ Предусмотрено & \centering -- & \centering \arraybackslash -- \\
\hline
\centering Порядок действия (Очередность\_ВНР)  & \centering SGF2 & \centering ВВ - СВ \\ СВ - ВВ & \centering -- & \centering \arraybackslash -- \\
\hline
\centering Время ввода ВНР (Tввода)  & \centering T1 & \centering 10000...170000 & \centering мс & \centering \arraybackslash 10 \\
\hline
\centering Выдержка времени 1 операции (Топер1)  & \centering T2 & \centering 0...50000 & \centering мс & \centering \arraybackslash 10 \\
\hline
\centering Выдержка времени 2 операции (Топер2)  & \centering T3 & \centering 0...50000 & \centering мс & \centering \arraybackslash 10 \\
\hline
\centering Время ожидания включения ввода (TвклВВ)  & \centering T4 & \centering 1000...65000 & \centering мс & \centering \arraybackslash 10 \\
\hline
%===t1
\end{longtable}
\normalsize


%%%%%%%%%%%%%%%%%%%%%%%%%%%%%%%%%%%%%%%%%%%%%%%%%%%%%%%%%% ВНР %%%%%%%%%%%%%%%%%%%%%%%%%%%%%%%%%%%%%%%%%%%%%%%%%%%%%%%%%%%%%%%%%%
\par
В таблице \ref{lorz:tbl1} приведены параметры, необходимые для настройки ЛО.


\footnotesize
\begin{longtable}{|>{\centering\arraybackslash}m{5.3cm}|>{\centering\arraybackslash}m{3.3cm}|>{\centering\arraybackslash}m{4.2cm}|>{\centering\arraybackslash}m{1.8cm}|>{\centering\arraybackslash}m{1cm}|}
\caption{Параметры для настройки ЛО РЗ\hfill\vspace{-0.5\baselineskip}}\label{lorz:tbl1}\\ 
\hline
\rowcolor{unimain!40}
Параметр (Параметр на ИЧМ) & Условное обозначение на схеме & Значение/ Диапазон & Единица измерения & Шаг \\ 
\hline
\endfirsthead
\caption*{\hspace{3pt}\emph{Продолжение таблицы \ref{lorz:tbl1}\hfill\vspace{-0.5\baselineskip}}} \\ % сделано по ГОСТ 2.105 п.6.8.7
\hline
\rowcolor{unimain!40}
Параметр (Параметр на ИЧМ) & Условное обозначение на схеме & Значение/ Диапазон & Единица измерения & Шаг \\ 
%\hline
\endhead
%\hline
\endfoot
%\hline
\endlastfoot
%==+t1*OFFPTRC1|LVOFFLGC
\multicolumn{5}{|c|}{ ЛО РЗ } \\ \hline 
\centering Ввод функции в работу (Ввод\_функции)  & \centering SGF1 & \centering Не предусмотрено \\ Предусмотрено & \centering -- & \centering \arraybackslash -- \\
\hline
\centering Длительность импульса (Tимп)  & \centering Т1 & \centering 50...60000 & \centering мс & \centering \arraybackslash 10 \\
\hline
\multicolumn{5}{|c|}{ ЛО ВНР } \\ \hline 
\centering Ввод функции в работу (Ввод\_функции)  & \centering SGF1 & \centering Не предусмотрено \\ Предусмотрено & \centering -- & \centering \arraybackslash -- \\
\hline
\centering Длительность импульса (Tимп)  & \centering Т1 & \centering 50...60000 & \centering мс & \centering \arraybackslash 10 \\
\hline
\multicolumn{5}{|c|}{ ЗАВР } \\ \hline 
\centering Режим запрета АВР при наличии отключения внешнего 1 (ЗапАВР\_Откл\_внеш\_1)  & \centering SGF1 & \centering Не предусмотрено \\ Предусмотрено & \centering -- & \centering \arraybackslash -- \\
\hline
\centering Режим запрета АВР при наличии отключения внешнего 2 (ЗапАВР\_Откл\_внеш\_2)  & \centering SGF2 & \centering Не предусмотрено \\ Предусмотрено & \centering -- & \centering \arraybackslash -- \\
\hline
\centering Режим запрета АВР при наличии отключения внешнего 3 (ЗапАВР\_Откл\_внеш\_3)  & \centering SGF3 & \centering Не предусмотрено \\ Предусмотрено & \centering -- & \centering \arraybackslash -- \\
\hline
\centering Режим запрета АВР при наличии отключения внешнего 4 (ЗапАВР\_Откл\_внеш\_4)  & \centering SGF4 & \centering Не предусмотрено \\ Предусмотрено & \centering -- & \centering \arraybackslash -- \\
\hline
%===t1
\end{longtable}
\normalsize





%%%%%%%%%%%%%%%%%%%%%%%%%%%%%%%%%%%%%%%%%%%%%%%%%%%%%%%%%% КП %%%%%%%%%%%%%%%%%%%%%%%%%%%%%%%%%%%%%%%%%%%%%%%%%%%%%%%%%%%%%%%%%% \footnotesize
\par
 В таблице \ref{kp:tbl_gen} приведены параметры, необходимые для настройки КП.

\footnotesize
\begin{longtable}{|>{\centering\arraybackslash}m{5.3cm}|>{\centering\arraybackslash}m{3.3cm}|>{\centering\arraybackslash}m{4.2cm}|>{\centering\arraybackslash}m{1.8cm}|>{\centering\arraybackslash}m{1cm}|}
\caption{Параметры для настройки КП\hfill\vspace{-0.5\baselineskip}}\label{kp:tbl_gen}\\ 
\hline
\rowcolor{gray!20}
Параметр (Параметр на ИЧМ) & Условное обозначение на схеме & Значение/ Диапазон & Единица измерения & Шаг \\ 
\hline
\endfirsthead
\caption*{\hspace{3pt}\emph{Продолжение таблицы \ref{kp:tbl_gen}\hfill\vspace{-0.5\baselineskip}}} \\ % сделано по ГОСТ 2.105 п.6.8.7
\hline
\rowcolor{gray!20}
Параметр (Параметр на ИЧМ) & Условное обозначение на схеме & Значение/ Диапазон & Единица измерения & Шаг \\ 
%\hline
\endhead
%\hline
\endfoot
%\hline
\endlastfoot
%==+t1*LLN0|SV_SWCTRL
\multicolumn{5}{|c|}{ УВ } \\ \hline 
\centering Ввод функции в работу (Ввод\_функции)  & \centering SGF1 & \centering Не предусмотрено \\ Предусмотрено & \centering -- & \centering \arraybackslash -- \\
\hline
\centering Включение с контролем синхронизма (Контроль\_синхр)  & \centering SGF2 & \centering Не предусмотрено \\ Предусмотрено & \centering -- & \centering \arraybackslash -- \\
\hline
\centering Блокировка включения от аварийного отключения В (Бл\_вкл\_от\_авар\_откл)  & \centering SGF3 & \centering Не предусмотрено \\ Предусмотрено & \centering -- & \centering \arraybackslash -- \\
\hline
\centering Допустимое время переключения выключателя (Tперекл)  & \centering T1 & \centering 10...10000 & \centering мс & \centering \arraybackslash 10 \\
\hline
\centering Выдержка времени подавления выдачи положения "Не определено" (Tблк)  & \centering T2 & \centering 10...10000 & \centering мс & \centering \arraybackslash 10 \\
\hline
\centering Длительность задержки включения с контролем синхронизма (Тзад\_вкл\_с\_КС)  & \centering T3 & \centering 50...60000 & \centering мс & \centering \arraybackslash 10 \\
\hline
\multicolumn{5}{|c|}{ УВЭ1 } \\ \hline 
\centering Ввод функции в работу (Ввод\_функции)  & \centering SGF1 & \centering Не предусмотрено \\ Предусмотрено & \centering -- & \centering \arraybackslash -- \\
\hline
\centering Допустимое время переключения (Tперекл)  & \centering T1 & \centering 50...100000 & \centering мс & \centering \arraybackslash 10 \\
\hline
\centering Выдержка времени подавления выдачи положения "Не определено" (Тблк)  & \centering T2 & \centering 50...100000 & \centering мс & \centering \arraybackslash 10 \\
\hline
\multicolumn{5}{|c|}{ УВЭ2 } \\ \hline 
\centering Ввод функции в работу (Ввод\_функции)  & \centering SGF1 & \centering Не предусмотрено \\ Предусмотрено & \centering -- & \centering \arraybackslash -- \\
\hline
\centering Допустимое время переключения (Tперекл)  & \centering T1 & \centering 50...100000 & \centering мс & \centering \arraybackslash 10 \\
\hline
\centering Выдержка времени подавления выдачи положения "Не определено" (Тблк)  & \centering T2 & \centering 50...100000 & \centering мс & \centering \arraybackslash 10 \\
\hline
\multicolumn{5}{|c|}{ УЗН } \\ \hline 
\centering Ввод функции в работу (Ввод\_функции)  & \centering SGF1 & \centering Не предусмотрено \\ Предусмотрено & \centering -- & \centering \arraybackslash -- \\
\hline
\centering Допустимое время переключения (Tперекл)  & \centering T1 & \centering 50...100000 & \centering мс & \centering \arraybackslash 10 \\
\hline
\centering Выдержка времени подавления выдачи положения "Не определено" (Тблк)  & \centering T2 & \centering 50...100000 & \centering мс & \centering \arraybackslash 10 \\
\hline
\multicolumn{5}{|c|}{ ОБР В } \\ \hline 
\centering Ввод ОБР В в работу (ОБР\_В)  & \centering SGF1 & \centering Не предусмотрено \\ Предусмотрено & \centering -- & \centering \arraybackslash -- \\
\hline
\centering Контроль ЗН (Контр\_ЗН)  & \centering SGF2 & \centering Не предусмотрено \\ Предусмотрено & \centering -- & \centering \arraybackslash -- \\
\hline
\centering Контроль КА1 (Контр\_KA1)  & \centering SGF3 & \centering Не предусмотрено \\ Предусмотрено & \centering -- & \centering \arraybackslash -- \\
\hline
\centering Контроль КА2 (Контр\_KA2)  & \centering SGF4 & \centering Не предусмотрено \\ Предусмотрено & \centering -- & \centering \arraybackslash -- \\
\hline
\centering Контроль от ВЭ2 (Контр\_от\_ВЭ2)  & \centering SGF5 & \centering Не предусмотрено \\ Предусмотрено & \centering -- & \centering \arraybackslash -- \\
\hline
\centering Контроль КА4 (Контр\_KA4)  & \centering SGF6 & \centering Не предусмотрено \\ Предусмотрено & \centering -- & \centering \arraybackslash -- \\
\hline
\multicolumn{5}{|c|}{ ОБР ВЭ1 } \\ \hline 
\centering Контроль КА1 (Контр\_КА1)  & \centering SGF1 & \centering Не предусмотрено \\ Предусмотрено & \centering -- & \centering \arraybackslash -- \\
\hline
\centering Контроль КА2 (Контр\_КА2)  & \centering SGF2 & \centering Не предусмотрено \\ Предусмотрено & \centering -- & \centering \arraybackslash -- \\
\hline
\centering Контроль КА4 (Контр\_КА4)  & \centering SGF3 & \centering Не предусмотрено \\ Предусмотрено & \centering -- & \centering \arraybackslash -- \\
\hline
\multicolumn{5}{|c|}{ ОБР ВЭ2 } \\ \hline 
\centering Контроль КА5 (Контр\_КА5)  & \centering SGF1 & \centering Не предусмотрено \\ Предусмотрено & \centering -- & \centering \arraybackslash -- \\
\hline
\multicolumn{5}{|c|}{ КП } \\ \hline 
\centering Ввод функции в работу (Ввод\_функции)  & \centering SGF1 & \centering Не предусмотрено \\ Предусмотрено & \centering -- & \centering \arraybackslash -- \\
\hline
\centering Ввод ОБР КА в работу (ОБР\_КА)  & \centering SG1 & \centering Не предусмотрено \\ Предусмотрено & \centering -- & \centering \arraybackslash -- \\
\hline
%===t1
\end{longtable}
\normalsize



%%%%%%%%%%%%%%%%%%%%%%%%%%%%%%%%%%%%%%%%%%%%%%%%%%%%%%%%%% УВ %%%%%%%%%%%%%%%%%%%%%%%%%%%%%%%%%%%%%%%%%%%%%%%%%%%%%%%%%%%%%%%%%%
\par
В таблице \ref{uv:tbl_gen} приведены параметры, необходимые для настройки УВ. 

\footnotesize
\begin{longtable}{|>{\centering\arraybackslash}m{5.3cm}|>{\centering\arraybackslash}m{3.3cm}|>{\centering\arraybackslash}m{4.2cm}|>{\centering\arraybackslash}m{1.8cm}|>{\centering\arraybackslash}m{1cm}|}
\caption{Параметры для настройки УВ\hfill\vspace{-0.5\baselineskip}}\label{uv:tbl_gen}\\ 
\hline
\rowcolor{gray!20}
Параметр (Параметр на ИЧМ) & Условное обозначение на схеме & Значение/ Диапазон & Единица измерения & Шаг \\ 
\hline
\endfirsthead
\caption*{\hspace{3pt}\emph{Продолжение таблицы \ref{uv:tbl_gen}\hfill\vspace{-0.5\baselineskip}}} \\ % сделано по ГОСТ 2.105 п.6.8.7
\hline
\rowcolor{gray!20}
Параметр (Параметр на ИЧМ) & Условное обозначение на схеме & Значение/ Диапазон & Единица измерения & Шаг \\ 
%\hline
\endhead
%\hline
\endfoot
%\hline
\endlastfoot
%==+t1*CBCSWI1|HVBCTRL
\centering Ввод функции в работу (Ввод\_функции)  & \centering SGF1 & \centering Не предусмотрено \\ Предусмотрено & \centering -- & \centering \arraybackslash -- \\
\hline
\centering Включение с контролем синхронизма (Контроль\_синхр)  & \centering SGF2 & \centering Не предусмотрено \\ Предусмотрено & \centering -- & \centering \arraybackslash -- \\
\hline
\centering Длительность задержки включения с контролем синхронизма (Тзад\_вкл\_с\_КС)  & \centering T1 & \centering 50...60000 & \centering мс & \centering \arraybackslash 10 \\
\hline
%===t1
\end{longtable}
\normalsize



%%%%%%%%%%%%%%%%%%%%%%%%%%%%%%%%%%%%%%%%%%%%%%%%%%%%%%%%%% КА %%%%%%%%%%%%%%%%%%%%%%%%%%%%%%%%%%%%%%%%%%%%%%%%%%%%%%%%%%%%%%%%%%
\par
В таблице \ref{ka:tbl_gen} приведены параметры, необходимые для настройки КА.

\footnotesize
\begin{longtable}{|>{\centering\arraybackslash}m{5.3cm}|>{\centering\arraybackslash}m{3.3cm}|>{\centering\arraybackslash}m{4.2cm}|>{\centering\arraybackslash}m{1.8cm}|>{\centering\arraybackslash}m{1cm}|}
\caption{Параметры для настройки КА\hfill\vspace{-0.5\baselineskip}}\label{ka:tbl_gen}\\ 
\hline
\rowcolor{gray!20}
Параметр (Параметр на ИЧМ) & Условное обозначение на схеме & Значение/ Диапазон & Единица измерения & Шаг \\ 
\hline
\endfirsthead
%\caption{\emph{Продолжение\hfill\vspace{-0.5\baselineskip}}} \\ % Отступ обозначения таблицы от таблицы и выравнивание надписи слева
\caption*{\hspace{3pt}\emph{Продолжение таблицы \ref{ka:tbl_gen}\hfill\vspace{-0.5\baselineskip}}} \\ % сделано по ГОСТ 2.105 п.6.8.7
\hline
\rowcolor{gray!20}
Параметр (Параметр на ИЧМ) & Условное обозначение на схеме & Значение/ Диапазон & Единица измерения & Шаг \\ 
%\hline
\endhead
%\hline
\endfoot
%\hline
\endlastfoot
%==+t1*LLN0|SV_SwitchDevice
\multicolumn{5}{|c|}{ В } \\ \hline 
\centering Ввод функции в работу (Ввод\_функции)  & \centering SGF1 & \centering Не предусмотрено \\ Предусмотрено & \centering -- & \centering \arraybackslash -- \\
\hline
\centering Режим отключения выключателя (Режим\_откл)  & \centering SGF2 & \centering Длительный \\ Импульсный & \centering -- & \centering \arraybackslash -- \\
\hline
\centering Контроль работы ЭМО (Контр\_раб\_ЭМО)  & \centering SGF3 & \centering Не предусмотрено \\ Предусмотрено & \centering -- & \centering \arraybackslash -- \\
\hline
\centering Режим включения выключателя (Режим\_вкл)  & \centering SGF4 & \centering Длительный \\ Импульсный & \centering -- & \centering \arraybackslash -- \\
\hline
\centering Контроль работы ЭМВ (Контр\_раб\_ЭМВ)  & \centering SGF5 & \centering Не предусмотрено \\ Предусмотрено & \centering -- & \centering \arraybackslash -- \\
\hline
\centering Длительность формирования сигнала о неисправном положении В (Tнеиспр\_В)  & \centering T1 & \centering 10...10000 & \centering мс & \centering \arraybackslash 10 \\
\hline
\centering Длительность формирования команды "Отключить В (реле)" (Тимп\_откл)  & \centering T2 & \centering 50...10000 & \centering мс & \centering \arraybackslash 10 \\
\hline
\centering Длительность формирования команды "Включить В (реле)" (Тимп\_вкл)  & \centering T3 & \centering 50...10000 & \centering мс & \centering \arraybackslash 10 \\
\hline
\centering Время для исключения "опрокидывания" В (Тпрод\_вкл)  & \centering T4 & \centering 0...1000 & \centering мс & \centering \arraybackslash 10 \\
\hline
\multicolumn{5}{|c|}{ ЗН } \\ \hline 
\centering Ввод функции в работу (Ввод\_функции)  & \centering SGF1 & \centering Не предусмотрено \\ Предусмотрено & \centering -- & \centering \arraybackslash -- \\
\hline
\centering Режим отключения ЗН (Режим\_откл)  & \centering SGF2 & \centering Длительный \\ Импульсный & \centering -- & \centering \arraybackslash -- \\
\hline
\centering Режим включения ЗН (Режим\_вкл)  & \centering SGF3 & \centering Длительный \\ Импульсный & \centering -- & \centering \arraybackslash -- \\
\hline
\centering Длительность формирования сигнала о неисправном положении ЗН (Tнеиспр\_ЗН)  & \centering T1 & \centering 10...100000 & \centering мс & \centering \arraybackslash 10 \\
\hline
\centering Длительность формирования команды "Отключить ЗН (реле)" (Тимп\_откл)  & \centering T2 & \centering 100...100000 & \centering мс & \centering \arraybackslash 10 \\
\hline
\centering Длительность формирования команды "Включить ЗН (реле)" (Тимп\_вкл)  & \centering T3 & \centering 100...100000 & \centering мс & \centering \arraybackslash 10 \\
\hline
\multicolumn{5}{|c|}{ ВЭ1 } \\ \hline 
\centering Ввод функции в работу (Ввод\_функции)  & \centering SGF1 & \centering Не предусмотрено \\ Предусмотрено & \centering -- & \centering \arraybackslash -- \\
\hline
\centering Режим работы (выкатить) ВЭ (Режим\_выкат)  & \centering SGF2 & \centering Длительный \\ Импульсный & \centering -- & \centering \arraybackslash -- \\
\hline
\centering Режим работы (вкатить) ВЭ (Режим\_вкат)  & \centering SGF3 & \centering Длительный \\ Импульсный & \centering -- & \centering \arraybackslash -- \\
\hline
\centering Контроль ремонтного положения (Контр\_рем\_полож)  & \centering SGF4 & \centering Не предусмотрено \\ Предусмотрено & \centering -- & \centering \arraybackslash -- \\
\hline
\centering Длительность формирования сигнала о неисправном положении ВЭ (Tнеиспр\_ВЭ)  & \centering T1 & \centering 10...100000 & \centering мс & \centering \arraybackslash 10 \\
\hline
\centering Длительность формирования команды "Выкатить ВЭ (реле)" (Тимп\_выкат)  & \centering T2 & \centering 100...100000 & \centering мс & \centering \arraybackslash 10 \\
\hline
\centering Длительность формирования команды "Вкатить ВЭ (реле)" (Тимп\_вкат)  & \centering T3 & \centering 100...100000 & \centering мс & \centering \arraybackslash 10 \\
\hline
\multicolumn{5}{|c|}{ ВЭ2 } \\ \hline 
\centering Ввод функции в работу (Ввод\_функции)  & \centering SGF1 & \centering Не предусмотрено \\ Предусмотрено & \centering -- & \centering \arraybackslash -- \\
\hline
\centering Режим работы (выкатить) ВЭ (Режим\_выкат)  & \centering SGF2 & \centering Длительный \\ Импульсный & \centering -- & \centering \arraybackslash -- \\
\hline
\centering Режим работы (вкатить) ВЭ (Режим\_вкат)  & \centering SGF3 & \centering Длительный \\ Импульсный & \centering -- & \centering \arraybackslash -- \\
\hline
\centering Контроль ремонтного положения (Контр\_рем\_полож)  & \centering SGF4 & \centering Не предусмотрено \\ Предусмотрено & \centering -- & \centering \arraybackslash -- \\
\hline
\centering Длительность формирования сигнала о неисправном положении ВЭ (Tнеиспр\_ВЭ)  & \centering T1 & \centering 10...100000 & \centering мс & \centering \arraybackslash 10 \\
\hline
\centering Длительность формирования команды "Выкатить ВЭ (реле)" (Тимп\_выкат)  & \centering T2 & \centering 100...100000 & \centering мс & \centering \arraybackslash 10 \\
\hline
\centering Длительность формирования команды "Вкатить ВЭ (реле)" (Тимп\_вкат)  & \centering T3 & \centering 100...100000 & \centering мс & \centering \arraybackslash 10 \\
\hline
\multicolumn{5}{|c|}{ КА } \\ \hline 
\centering Ввод функции в работу (Ввод\_функции)  & \centering SGF1 & \centering Не предусмотрено \\ Предусмотрено & \centering -- & \centering \arraybackslash -- \\
\hline
%===t1
\end{longtable}
\normalsize



%%%%%%%%%%%%%%%%%%%%%%%%%%%%%%%%%%%%%%%%%%%%%%%%%%%%%%%%%% КСВ %%%%%%%%%%%%%%%%%%%%%%%%%%%%%%%%%%%%%%%%%%%%%%%%%%%%%%%%%%%%%%%%%%
\par
В таблице \ref{ksv:tbl1} приведены параметры, необходимые для настройки КСВ.

\footnotesize
\begin{longtable}{|>{\centering\arraybackslash}m{5.3cm}|>{\centering\arraybackslash}m{3.3cm}|>{\centering\arraybackslash}m{4.2cm}|>{\centering\arraybackslash}m{1.8cm}|>{\centering\arraybackslash}m{1cm}|}
\caption{Параметры для настройки КСВ\hfill\vspace{-0.5\baselineskip}}\label{ksv:tbl1}\\ 
\hline
\rowcolor{gray!20}
Параметр (Параметр на ИЧМ) & Условное обозначение на схеме & Значение/ Диапазон & Единица измерения & Шаг \\ 
\hline
\endfirsthead
\caption*{\hspace{3pt}\emph{Продолжение таблицы \ref{ksv:tbl1}\hfill\vspace{-0.5\baselineskip}}} \\ % сделано по ГОСТ 2.105 п.6.8.7
\hline
\rowcolor{gray!20}
Параметр (Параметр на ИЧМ) & Условное обозначение на схеме & Значение/ Диапазон & Единица измерения & Шаг \\ 
%\hline
\endhead
%\hline
\endfoot
%\hline
\endlastfoot
%==+t1*RCBF1|LVCBSUP
\centering Ввод функции в работу (Ввод\_функции)  & \centering SGF1 & \centering Не предусмотрено \\ Предусмотрено & \centering -- & \centering \arraybackslash -- \\
\hline
\centering Контроль ОТ цепей ЭМВ, ЭМО1 и ЭМО2 (Контроль\_ОТ\_ЭМ)  & \centering SGF5 & \centering Не предусмотрено \\ ЭМВ и ЭМО1 \\ ЭМВ, ЭМО1 и ЭМО2 & \centering -- & \centering \arraybackslash -- \\
\hline
\centering Блокировка управления при снижении уровня изоляции В (Блок\_упр\_КИ\_В)  & \centering SGF8 & \centering Не предусмотрено \\ От аварийного \\ От аварийного и низкого & \centering -- & \centering \arraybackslash -- \\
\hline
\centering Блокировка включения при низком уровне изоляции В (Блок\_вкл\_низ\_изол\_В)  & \centering SGF2 & \centering Не предусмотрено \\ Предусмотрено & \centering -- & \centering \arraybackslash -- \\
\hline
\centering Блокировка включения при неисправности положения В (Блок\_вкл\_полож\_В)  & \centering SGF3 & \centering Не предусмотрено \\ Предусмотрено & \centering -- & \centering \arraybackslash -- \\
\hline
\centering Блокировка включения при превышении ресурса В (Блок\_вкл\_ресурса\_В)  & \centering SGF4 & \centering Не предусмотрено \\ Предусмотрено & \centering -- & \centering \arraybackslash -- \\
\hline
\centering Контроль ЭМВ, ЭМО1 и ЭМО2 при формировании неисправности цепей ЭМУ (Контроль\_ЭМ)  & \centering SGF6 & \centering Не предусмотрено \\ ЭМВ и ЭМО1 \\ ЭМВ, ЭМО1 и ЭМО2 & \centering -- & \centering \arraybackslash -- \\
\hline
\centering Разрешение сброса "РФК" от кнопки (Контроль\_кнопки)  & \centering SGF7 & \centering Не предусмотрено \\ Предусмотрено & \centering -- & \centering \arraybackslash -- \\
\hline
\centering Выдержка времени ожидания готовности привода выключателя (Tгот\_в)  & \centering T1 & \centering 0...100000 & \centering мс & \centering \arraybackslash 10 \\
\hline
\centering Выдержка времени срабатывания неисправности цепей ЭМУ (Tдлит\_неисп\_ЭМУ)  & \centering T2 & \centering 0...10000 & \centering мс & \centering \arraybackslash 10 \\
\hline
\centering Выдержка времени срабатывания защит электромагнитов (Tдлит\_ЭМ)  & \centering T3 & \centering 0...10000 & \centering мс & \centering \arraybackslash 10 \\
\hline
%===t1
\end{longtable}
\normalsize



%%%%%%%%%%%%%%%%%%%%%%%%%%%%%%%%%%%%%%%%%%%%%%%%%%%%%%%%%% КРВ %%%%%%%%%%%%%%%%%%%%%%%%%%%%%%%%%%%%%%%%%%%%%%%%%%%%%%%%%%%%%%%%%%
\par
В таблице \ref{krv:tbl1} приведены параметры, необходимые для настройки КРВ.

\footnotesize
\begin{longtable}{|>{\centering\arraybackslash}m{5.3cm}|>{\centering\arraybackslash}m{3.3cm}|>{\centering\arraybackslash}m{4.2cm}|>{\centering\arraybackslash}m{1.8cm}|>{\centering\arraybackslash}m{1cm}|}
\caption{Параметры для настройки КРВ\hfill\vspace{-0.5\baselineskip}}\label{krv:tbl1}\\ 
\hline
\rowcolor{gray!20} Параметр (Параметр на ИЧМ) & Условное обозначение на схеме & Значение/ Диапазон & Единица измерения & Шаг \\ 
\hline
\endfirsthead
\caption*{\hspace{3pt}\emph{Продолжение таблицы \ref{krv:tbl1}\hfill\vspace{-0.5\baselineskip}}} \\ % сделано по ГОСТ 2.105 п.6.8.7
\hline
\rowcolor{gray!20} Параметр (Параметр на ИЧМ) & Условное обозначение на схеме & Значение/ Диапазон & Единица измерения & Шаг \\ 
%\hline
\endhead
%\hline
\endfoot
%\hline
\endlastfoot 
%==+t1*CLS|CLS
\centering Ввод функции в работу (Ввод\_функции)  & \centering SGF1 & \centering Не предусмотрено \\ Предусмотрено & \centering -- & \centering \arraybackslash -- \\
\hline
\centering Минимальный ток КРВ (Iмин\_КРВ)  & \centering -- & \centering 0,01...0,30 & \centering о.е. & \centering \arraybackslash 0,01 \\
\hline
\centering Номинальный ток выключателя (паспорт) (Iном\_В\_пасп)  & \centering -- & \centering 0...5000 & \centering А & \centering \arraybackslash 1 \\
\hline
\centering Номинальный ток отключения выключателя (паспорт) (Iном\_откл\_В\_пасп)  & \centering -- & \centering 0...40000 & \centering А & \centering \arraybackslash 1 \\
\hline
\centering Допустимое количество циклов при номинальном токе выключателя (паспорт) (КРВпасп\_Iном)  & \centering -- & \centering 0...200000 & \centering -- & \centering \arraybackslash 1 \\
\hline
\centering Допустимое количество циклов при номинальном токе отключения выключателя (паспорт) (КРВпасп\_Iном\_откл)  & \centering -- & \centering 0...500 & \centering -- & \centering \arraybackslash 1 \\
\hline
\centering Допустимый механический ресурс выключателя (паспорт) (МРВпасп)  & \centering -- & \centering 0...200000 & \centering -- & \centering \arraybackslash 1 \\
\hline
\centering Начальное значение коммутационного ресурса выключателя (Нач\_знач\_КРВ)  & \centering -- & \centering 0,0...100,0 & \centering \% & \centering \arraybackslash 0,1 \\
\hline
\centering Величина срабатывания коммутационного ресурса выключателя (КРВсраб)  & \centering -- & \centering 0,0...100,0 & \centering \% & \centering \arraybackslash 0,1 \\
\hline
\centering Начальное значение механического ресурса выключателя (Нач\_знач\_МРВ)  & \centering -- & \centering 0...200000 & \centering -- & \centering \arraybackslash 1 \\
\hline
\centering Контроль КРВ по превышению МРВ (Контроль\_КРВ\_от\_МРВ)  & \centering SGF2 & \centering Не предусмотрено \\ Предусмотрено & \centering -- & \centering \arraybackslash -- \\
\hline
\centering Полное время отключения выключателя (Tмакс\_В)  & \centering -- & \centering 10...3000 & \centering мс & \centering \arraybackslash 10 \\
\hline
%===t1
\end{longtable}
\normalsize


%%%%%%%%%%%%%%%%%%%%%%%%%%%%%%%%%%%%%%%%%%%%%%%%%%%%%%%%%% ПС %%%%%%%%%%%%%%%%%%%%%%%%%%%%%%%%%%%%%%%%%%%%%%%%%%%%%%%%%%%%%%%%%%
\par
В таблице \ref{ps:tbl1} приведены параметры, необходимые для настройки ПС. 

\footnotesize
\begin{longtable}{|>{\centering\arraybackslash}m{5.3cm}|>{\centering\arraybackslash}m{3.3cm}|>{\centering\arraybackslash}m{4.2cm}|>{\centering\arraybackslash}m{1.8cm}|>{\centering\arraybackslash}m{1cm}|}
\caption{Параметры для настройки предупредительной сигнализации\hfill\vspace{-0.5\baselineskip}}\label{ps:tbl1}\\ 
\hline
\rowcolor{unimain!40}
Параметр (Параметр на ИЧМ) & Условное обозначение на схеме & Значение/ Диапазон & Единица измерения & Шаг \\ 
\hline
\endfirsthead
%\caption{\emph{Продолжение\hfill\vspace{-0.5\baselineskip}}} \\ % Отступ обозначения таблицы от таблицы и выравнивание надписи слева
\caption*{\hspace{3pt}\emph{Продолжение таблицы \ref{ps:tbl1}\hfill\vspace{-0.5\baselineskip}}} \\ % сделано по ГОСТ 2.105 п.6.8.7
\hline
\rowcolor{unimain!40}
Параметр (Параметр на ИЧМ) & Условное обозначение на схеме & Значение/ Диапазон & Единица измерения & Шаг \\ 
%\hline
\endhead
%\hline
\endfoot
%\hline
\endlastfoot
%==+t1*CALH1|SV_LVALH
\centering Контроль выходных цепей (Контр\_цепей)  & \centering SGF1 & \centering Не предусмотрено \\ Предусмотрено & \centering -- & \centering \arraybackslash -- \\
\hline
\centering Контроль выведенного положения БИ (Контр\_вывода\_БИ)  & \centering SGF2 & \centering Не предусмотрено \\ Предусмотрено & \centering -- & \centering \arraybackslash -- \\
\hline
\centering Контроль сигнализации опер.тока (Контр\_ОТ\_сигн)  & \centering SGF3 & \centering Не предусмотрено \\ Предусмотрено & \centering -- & \centering \arraybackslash -- \\
\hline
\centering Контроль неисправности КА (Контр\_Неиспр\_КА)  & \centering SGF4 & \centering Не предусмотрено \\ Предусмотрено & \centering -- & \centering \arraybackslash -- \\
\hline
\centering Контроль превышения времени переключения КА (Контр\_прев\_вр\_пер)  & \centering SGF5 & \centering Не предусмотрено \\ Предусмотрено & \centering -- & \centering \arraybackslash -- \\
\hline
\centering Контроль общего внешнего сигнала (Контр\_общ\_вн\_сигн)  & \centering SGF6 & \centering Не предусмотрено \\ Предусмотрено & \centering -- & \centering \arraybackslash -- \\
\hline
\centering Контроль положения переключателя 1 аварийной деблокировки (Контр\_пол\_ав\_дебл\_1)  & \centering SGF7 & \centering Не предусмотрено \\ Предусмотрено & \centering -- & \centering \arraybackslash -- \\
\hline
\centering Контроль положения переключателя 2 аварийной деблокировки (Контр\_пол\_ав\_дебл\_2)  & \centering SGF8 & \centering Не предусмотрено \\ Предусмотрено & \centering -- & \centering \arraybackslash -- \\
\hline
%===t1
\end{longtable}
\normalsize


%%%%%%%%%%%%%%%%%%%%%%%%%%%%%%%%%%%%%%%%%%%%%%%%%%%%%%%%%% СС %%%%%%%%%%%%%%%%%%%%%%%%%%%%%%%%%%%%%%%%%%%%%%%%%%%%%%%%%%%%%%%%%%
\par
В таблице \ref{ss:tbl1} приведены параметры, необходимые для настройки СС. 

\footnotesize
\begin{longtable}{|>{\centering\arraybackslash}m{5.3cm}|>{\centering\arraybackslash}m{3.3cm}|>{\centering\arraybackslash}m{4.2cm}|>{\centering\arraybackslash}m{1.8cm}|>{\centering\arraybackslash}m{1cm}|}
\caption{Параметры для настройки функции сборки сигналов\hfill\vspace{-0.5\baselineskip}}\label{ss:tbl1}\\ 
\hline
\rowcolor{unimain!40}
Параметр (Параметр на ИЧМ) & Условное обозначение на схеме & Значение/ Диапазон & Единица измерения & Шаг \\ 
\hline
\endfirsthead
\caption*{\hspace{3pt}\emph{Продолжение таблицы \ref{ss:tbl1}\hfill\vspace{-0.5\baselineskip}}} \\ % сделано по ГОСТ 2.105 п.6.8.7
\hline
\rowcolor{unimain!40}
Параметр (Параметр на ИЧМ) & Условное обозначение на схеме & Значение/ Диапазон & Единица измерения & Шаг \\ 
%\hline
\endhead
%\hline
\endfoot
%\hline
\endlastfoot
%==+t1*GAPC1|SV_SignAssembly
\multicolumn{5}{|c|}{ CC } \\ \hline 
\centering Контроль опер. тока цепей сигнализации выключателя (Контр\_ОТ\_сигн\_В)  & \centering SGF1 & \centering Не предусмотрено \\ Предусмотрено & \centering -- & \centering \arraybackslash -- \\
\hline
\centering Контроль опер. тока цепей ОБР 1 (Контр\_ОТ\_ОБР1)  & \centering SGF2 & \centering Не предусмотрено \\ Предусмотрено & \centering -- & \centering \arraybackslash -- \\
\hline
\centering Контроль опер. тока цепей ОБР 2 (Контр\_ОТ\_ОБР2)  & \centering SGF3 & \centering Не предусмотрено \\ Предусмотрено & \centering -- & \centering \arraybackslash -- \\
\hline
\centering Контроль положения SA1 (Контр\_полож\_SA1)  & \centering SGF4 & \centering Не предусмотрено \\ Предусмотрено & \centering -- & \centering \arraybackslash -- \\
\hline
\centering Контроль положения SA2 (Контр\_полож\_SA2)  & \centering SGF5 & \centering Не предусмотрено \\ Предусмотрено & \centering -- & \centering \arraybackslash -- \\
\hline
\centering Контроль положения SA3 (Контр\_полож\_SA3)  & \centering SGF6 & \centering Не предусмотрено \\ Предусмотрено & \centering -- & \centering \arraybackslash -- \\
\hline
\centering Контроль положения SA4 (Контр\_полож\_SA4)  & \centering SGF7 & \centering Не предусмотрено \\ Предусмотрено & \centering -- & \centering \arraybackslash -- \\
\hline
\centering Контроль положения SG1 (Контр\_полож\_SG1)  & \centering SGF8 & \centering Не предусмотрено \\ Предусмотрено & \centering -- & \centering \arraybackslash -- \\
\hline
\centering Контроль положения SG2 (Контр\_полож\_SG2)  & \centering SGF9 & \centering Не предусмотрено \\ Предусмотрено & \centering -- & \centering \arraybackslash -- \\
\hline
\centering Контроль положения SG3 (Контр\_полож\_SG3)  & \centering SGF10 & \centering Не предусмотрено \\ Предусмотрено & \centering -- & \centering \arraybackslash -- \\
\hline
%===t1
\end{longtable}
\normalsize